% The appendices of the notes: Appendix A is the method; then one chapter
% per backed claim.

\chapter{Checking a claim}
\label{app:method}

Most of what this report says comes from the session's own records: the
exported banks, the agents' change logs, the reviews and the builds.  Those
statements are about one case, and the report does not generalize them:
the lessons (\cref{sec:lessons}) and the advice for your own course
(\cref{sec:own-course}) are this case's experience, stated as such, and
no literature was searched for them.
A few statements are of another kind: advice about writing exam items that
rests on what is known beyond this case.  They cannot be checked at the
keyboard, and each has an appendix of its own.
\begin{itemize}
  \item Cues to the key and options that refer to other options are
    item-writing flaws, and such flaws matter
    (\cref{sec:items}; \cref{app:item-flaws}).
  \item Ordering items (Parsons problems) are used to assess programming,
    and an ordering item can have more than one valid order
    (\cref{sec:items,sec:checker}; \cref{app:parsons}).
  \item Feedback that explains supports learning more than feedback that
    only says right or wrong, though not always (\cref{sec:workflow};
    \cref{app:feedback}).
\end{itemize}
This appendix describes how those three claims were checked.

\section{Turn the claim into a question}

A claim can be checked only once it is clear what would show it wrong.
\enquote{Avoid \enquote{none of the above}} is advice; the question behind
it is whether items with that option are worse, by some measure, than items
without it.  Each appendix opens with such a question and says how strong
an answer the report needs.

\section{Choose the right kind of source}

A guideline's existence is shown by the guideline itself; whether
following it helps is shown by studies that compare items that follow it
with items that do not, and whether that is established by reviews and
meta-analyses that weigh many such studies.  Each appendix says which kind
of source answers which part of its claim.

\section{Search so that you can be wrong}

The searches were driven by the topic, not by authors: the queries name
the concepts of the claim, and a source enters because a query found it,
not because we expected it.  For every claim we also ran counter-queries,
which use the same concepts together with words for criticism and
limitations, so that evidence against the claim had the same chance to
surface as evidence for it.  The same queries were run in every database
that answered on the day, and every count of hits is reported, including
the failed searches.

\section{Read in full text, and write down how}

A hit is not evidence.  Every source the report cites was read in full
text where we could obtain it, and the passage that bears on the claim is
quoted with its page.  Where only the abstract could be read, the appendix
says so.  For every source the bibliography records why it was picked,
how it was found, the quoted passage, how it was verified, and what the
counter-search found.

\section{Draw the conclusion, with its reservations}

Each appendix ends by answering its question at the strength the sources
support, and by saying what the search left undone and what the field
would have to study for a surer answer.

\section{A note on language models}

The searches, the screening and the first reading were done by AI agents
with the search tool \texttt{scholar}, which records every query and its
results; a language model classified the hits as supporting, qualifying,
adjacent or off-topic, and the agents read every cited source themselves,
and the qualifying hits to the extent each appendix states.  A language
model can misclassify a hit or misread a passage, which is why every
quotation is given verbatim with its page, so that a reader can check it.


% ---------------------------------------------------------------- claim B
% Author's notes: scholar session "rapport-item-flaws"; export
%   scholar sessions export rapport-item-flaws -f table --bearing-only \
%     --lang en --label sok-item-flaws -o item-flaws-full
% (caption sentence on incomplete queries corrected by hand: it counted six
% combined-provider runs that the per-provider runs superseded).  Report,
% CSV and verbatim queries in the agent's scratch directory of 2026-11-13;
% provenance per source in ltnotes.bib.
\chapter{Are cues to the key flaws, and do flaws matter?}
\label{app:item-flaws}
\chapterprecis{The author has not yet reviewed the results in this
  appendix in full.}

\Cref{sec:items} calls two kinds of option a flaw: a key that stands out
from the other options, for instance by being the longest, and an option
that refers to the other options, such as \enquote{none of the above}.
\Cref{sec:exam} adds that a flawed item weighs heavily.  Behind these
statements are three questions.
\begin{description}
  \item[(a)] Do the established guidelines for writing multiple-choice
    items advise against cues to the key and against options that refer to
    other options?
  \item[(b)] Do such flaws measurably change how items behave, or who
    passes?
  \item[(c)] Do examinees use cues such as a longer key?
\end{description}
The first is a question about what the guidelines say, with the
counter-question whether they rest on evidence; the other two are
empirical, and the report needs them only at the strength
\enquote{measurably, in some studies}.  The report's own reason to avoid
options that refer to others, that the platform shuffles the options, is
not a question for the literature: it follows from the platform.  We
describe the search (\cref{sec:item-flaws-method}), what the sources say
(\cref{sec:item-flaws-results}), its limitations, the answer, and what is
left to do; the hit list closes the chapter.

\section{Method}
\label{sec:item-flaws-method}

The searches ran on 13 and 14 November 2026 with the tool
\texttt{scholar}.  Web of Science, Scopus, IEEE Xplore and OpenAlex
answered; Semantic Scholar refused our key and DBLP answered with a page
against automated clients, so neither was searched.  Every query ran
without a limit on the number of hits and with the same keywords in every
database: Scopus received them inside \texttt{TITLE-ABS-KEY(\dots)},
which searches titles, abstracts and keywords, as Web of Science's topic
search does; OpenAlex searches full text, which explains its larger
counts.  OpenAlex's daily budget, shared with the searches of
\cref{app:parsons,app:feedback}, ran out during the first round, and its
three counter-queries ran after the budget was reset.

Three support queries look for the guidelines (S1), for flaws together
with item statistics (S2), and for test-wiseness together with cues (S3).
Three counter-queries reuse their concept terms: with words for missing
evidence and criticism (M1), with the field's names for options that refer
to others together with item statistics (M2), and with words for null
effects (M3).  \Cref{tab:sok-queries-item-flaws} gives them with their
hits.

\begin{table}[htbp]
  \caption{Queries and hits per database for the claims of
    \cref{app:item-flaws}, 13--14 November 2026, tool \texttt{scholar}, no
    limit on hits.  MC stands for (\enquote{multiple-choice} OR
    \enquote{multiple choice}); quoted words were searched as phrases.
    S = support query, M = counter-query.  Databases: Web of Science
    (WoS), Scopus, IEEE Xplore (IEEE), OpenAlex (OA, which searches full
    text; not \enquote{open access}).  A 0 is an answer with no hits; no
    search failed in the end.}
  \label{tab:sok-queries-item-flaws}
  \centering\footnotesize
  \begin{tabularx}{\linewidth}{@{}l>{\raggedright\arraybackslash}X
    rrrr@{}}
    \toprule
       & Keywords & WoS & Scopus & IEEE & OA \\
    \midrule
    S1 & MC AND (\enquote{item-writing guidelines} OR \enquote{item writing
         guidelines} OR \enquote{item-writing rules} OR \enquote{item
         writing rules} OR \enquote{item-writing principles})
       & 51 & 53 & 0 & 1\,183 \\
    S2 & MC AND (\enquote{item-writing flaws} OR \enquote{item writing
         flaws} OR \enquote{flawed items} OR \enquote{item flaws} OR
         \enquote{item-writing violations}) AND (difficulty OR
         discrimination OR \enquote{pass rate} OR \enquote{pass-fail} OR
         \enquote{item analysis} OR \enquote{item statistics})
       & 45 & 40 & 0 & 601 \\
    S3 & MC AND (testwiseness OR \enquote{test-wiseness} OR \enquote{test
         wiseness} OR \enquote{test-wise} OR \enquote{test wise}) AND (cue
         OR cues OR cueing OR \enquote{option length} OR \enquote{longest
         option} OR \enquote{longest answer})
       & 3 & 5 & 0 & 360 \\
    \midrule
    M1 & the terms of S1 AND (empirical OR evidence OR critique OR
         criticism OR \enquote{lack of support})
       & 20 & 20 & 0 & 906 \\
    M2 & MC AND (\enquote{none of the above} OR \enquote{all of the above})
         AND (\enquote{item difficulty} OR \enquote{item discrimination} OR
         \enquote{item analysis} OR \enquote{item statistics} OR
         \enquote{item writing} OR \enquote{item-writing} OR \enquote{item
         format})
       & 17 & 14 & 0 & 1\,151 \\
    M3 & MC AND (the flaw terms of S2 OR the test-wiseness terms of S3 OR
         \enquote{option length} OR \enquote{longest option} OR
         \enquote{longest answer}) AND (negligible OR \enquote{no effect}
         OR \enquote{no significant} OR \enquote{little effect} OR
         \enquote{small effect} OR \enquote{did not affect} OR unaffected)
       & 11 & 11 & 0 & 1\,745 \\
    \bottomrule
  \end{tabularx}
\end{table}

The searches gave 3\,549 unique records.  A language model classified
every one of them as supporting, qualifying, adjacent or off-topic, and
gave each discarded record its reason: another field, another item format,
multiple-choice items used only as a research instrument, or item
generation without evidence on flaws.  We read every qualifying record and
the titles of the 278 supporting records the model was least sure of, and
corrected 18 classifications.  The model counted as supporting many
studies of how common flaws are, which do not test their effects, so the
supporting count in \cref{tab:sok-item-flaws} is an upper bound.  Twelve
of the 14 cited sources were read in full text.  Two, the 2002 review of
guidelines and a 2025 study of police selection tests, were read only in
their abstracts, since the publisher's site blocked automated access; the
2002 guidelines themselves were read in full as restated by one of their
authors in 2011.

\section{Results}
\label{sec:item-flaws-results}

\paragraph{What the guidelines say.} The standard reference is a taxonomy
of 31 guidelines that \textcquote{HaladynaDowningRodriguez2002}{was
validated through a logical process that included two sources of evidence:
the consensus achieved from reviewing what was found in 27 textbooks on
educational testing and the results of 27 research studies and reviews
published since 1990}.  One of its authors restates the guidelines on the
options in full.  Three of them read
\textcquote[205--206]{Rodriguez2011}{24. Keep the length of choices about
equal. 25. None-of-the-above should be used carefully. 26. Avoid
All-of-the-above.}  Two more ask to avoid clues to the right answer ---
grammatical inconsistencies, a conspicuous correct choice, pairs of
options that give the key away --- and to make all distractors
plausible\autocite[205--206]{Rodriguez2011}.  A review of 2025 gives the
same advice: options \textcquote[230]{Schmauder2025}{should be similar in
content, grammatical structure, and length}, and \enquote{all of the
above} and \enquote{none of the above}
\textcquote[230]{Schmauder2025}{should generally be avoided}.  Note that
the canonical list does not ban \enquote{none of the above}; it asks for
care.

\paragraph{Do the guidelines rest on evidence?} Mostly not, say four
sources independently.  Rodriguez writes that
\textcquote[204]{Rodriguez2011}{\textins*{m}ost of these guidelines were
based on logical reasoning and good writing practices; very few were based
on empirical evidence}, and that of the 14 guidelines on the options
\textcquote[205]{Rodriguez2011}{only 5 have been studied empirically}.  A
study of test-wiseness reports that for
\textcquote[020138-1]{DeVore2016}{nearly half of the rules no supporting
research could be identified}; a study of option guidelines calls their
importance \textcquote[326]{Martinez2009}{an assumption based more on
common sense than on the findings of well-founded research}; and a study
of \enquote{none of the above} notes that the guides' suggestions
\textcquote[1]{Pachai2015}{are infrequently based on empirical research}.
None of them says the guidelines are wrong; they say most are untested.

\paragraph{Do flaws change items and outcomes?}
In four medical-school exams, items that broke the guidelines
\textcquote[133]{Downing2005}{were 0--15 percentage points more difficult
than standard items measuring the same construct}, and the author
concludes that \textcquote[141]{Downing2005}{some students -- perhaps as
high as 10--15\% of students tested -- were incorrectly classified as
failed}.  His flaws were mostly of the kind that make items harder
(negative stems, \enquote{all} or \enquote{none of the above}, unfocused
stems), not cues, and he warns that the scales compared differ in length
and reliability.  In 7\,126 items from STEM courses, the flaws most
negatively associated with item quality included \enquote{longest option
correct} and \enquote{all of the above}, which
\textcquote[24]{SchmuckerMoore2026}{decreased the difficulty of the
question, while also weakening its discrimination}.  In 1\,925 veterinary
items, implausible distractors made items easier (95 against 80 per cent
correct) and less discriminating\autocite[6]{Rush2016}.

The same studies, and others, qualify this.  Rodriguez summarises a
meta-analysis: \textcquote[209]{Rodriguez2011}{The empirical evidence of
item-format effects was inconclusive.  When rule violations resulted in
statistically significant effects, they tended to be small.}  In 800
anatomy items the flaws \textcquote[1]{Pais2016}{explained 2.8 and 0.8\,\%}
of difficulty and discrimination, against 6.2 and 3.2 per cent for the
anatomical site.  In the veterinary items, \textcquote[9]{Rush2016}{many
item-writing flaws identified in this study did not impact the indices of
difficulty or discrimination value}; the longest-option cue, though common
there, is not among the flaws that did.
And in police selection tests the effects of poorly written items
\textcquote{Arthur2025}{were generally small}, though test takers
\enquote{consistently performed worse} on them.

\paragraph{None and all of the above.}
Here the counter-search found real counter-evidence.  In a classroom
experiment, \textcquote{Pachai2015}{NOTA increased question difficulty only
when it was the correct response, with no effect on difficulty of replacing
any distractor type with NOTA}, though it lowered discrimination when it
replaced a good distractor; the authors still advise against it.  In
another study, \enquote{all of the above} as the key even improved
discrimination\autocite[326]{Martinez2009}.  A formal analysis concludes
that \textcquote[271]{Persson2025}{there seems to be no harm in using this
format}.  Students do tend to avoid \enquote{none of the above} as a
distractor\autocite{DeVore2016}, which makes it a weak one.

\paragraph{Do examinees use cues?} In three experiments with a published
test for medical education, in which six of the eight cued items had a key
at least a third of a line longer than the other
options\autocite[252]{Willing2015}, \textcquote[247]{%
Willing2015}{\textins*{i}tems were answered correctly more often in the MC
condition (71\,\%) than in the DOMC condition (47\,\%), but only when
items contained cues to their solution}; in the second condition the
options are shown one at a time, so they cannot be compared.  The
meta-analysis summarised by Rodriguez found that
\textcquote[209]{Rodriguez2011}{\textins*{m}aking the correct option
longer (a common error made by classroom teachers) made items easier}.
Against this, a study of 23 items with a length violation found no such
effect\autocite[329]{Martinez2009}; a small experiment found flawed items
\textcquote[47]{Casu2019}{tended to be easier}, but the differences
\textcquote[47]{Casu2019}{were not systematic}; and test-wiseness effects
other than the avoidance of \enquote{none of the above} were
\textcquote[020138-8]{DeVore2016}{weak effects postinstruction}, for cues
other than length.

\section{Discussion and limitations}
\label{sec:item-flaws-discussion}

Semantic Scholar and DBLP could not be searched, and IEEE Xplore, which
indexes little educational measurement, returned nothing for any query.
Scopus returned no abstracts, so about 250 records were classified on
their titles alone, and the supporting records the model was less sure of
were checked on their titles, not their abstracts.  Two sources were read
only in abstract, and the meta-analysis behind Rodriguez's summary only
through that summary.  Most primary studies come from medical and other
professional education, most are observational, and the two experiments
on cues are small or use a single test.  We found no study of flaws in
programming exams, and none of computer-based exams that shuffle their
options.

\section{Conclusion}
\label{sec:item-flaws-conclusion}

Yes to (a), partly to (b), probably to (c).  The established guidelines
advise options of about equal length, no clues to the key, plausible
distractors, no \enquote{all of the above}, and care with \enquote{none of
the above}; but most of these rules rest on consensus rather than tests.
Where flaws have been studied they have changed difficulty and
discrimination, and in one study who passed, but across studies the
effects are usually small and vary with the flaw and the exam.  Examinees
do use cues such as a longer key in experiments, though other studies of
the length cue found weak or no effects.  For the report this means: the
rewrite removes cues because they reward something other than the
objective and because they can be measured in a bank, not because the
literature shows large effects; and it avoids options that refer to
others because the platform shuffles them, with the guidelines' advice
as a weaker, partly contested second reason.

\section{Further work}
\label{sec:item-flaws-further}

This search left the 2002 review and the 2025 police study to be read in
full (both closed behind a block for automated clients), the 1997
meta-analysis to be obtained, and several studies of \enquote{none of the
above} in the hit list unread, among them a meta-analysis from
1992.\footnote{Knowles and Welch, \emph{Educational and Psychological
Measurement}, \texttt{doi:10.1177/0013164492052003006}.}  Semantic Scholar
and DBLP should be searched when they answer.

The field lacks adequately powered experiments that vary one flaw at a
time, the length cue in particular, in real graded exams; any study of
flaws in programming courses; and any study of computer-based exams with
shuffled options.  If such a study found a large advantage from a longer
key, the report could state (c) without reservation; if it found none,
equal option lengths would rest on fairness and clarity rather than on
measured effects.

\section{Hits bearing on the claim}
\label{sec:item-flaws-hits}

\Cref{tab:sok-item-flaws} lists every record the screening kept as cited,
supporting or qualifying.
% --- scholar LaTeX fragment --------------------------------------------
% Generated by: scholar sessions export -f table
% Session: rapport-item-flaws
% \input this file from the body of a document whose
% preamble loads:
%   \usepackage{longtable}
%   \usepackage{booktabs}
%   \usepackage{array}
% ----------------------------------------------------------------------

\begingroup\footnotesize
\begin{longtable}{@{}%
  >{\raggedright\arraybackslash}p{0.44\textwidth}c%
  >{\raggedright\arraybackslash}p{0.13\textwidth}%
  >{\raggedright\arraybackslash}p{0.26\textwidth}@{}}
\caption{Records bearing on the claim, Item-writing flaws (14 cited, 353 supporting and 34 qualifying or contradicting it, out of 3549 unique hits; the 1222 adjacent and 1926 false hits are not listed). The last column gives the reason for each inclusion; for machine-classified rows the model's confidence is given. Databases: OpenAlex (OA; not open access), IEEE Xplore (IEEE), Web of Science (WoS), Scopus; the hits per query and database are in the query table of this appendix.}\label{tab:sok-item-flaws}\\
\toprule
Title & Year & Provider & Reason \\
\midrule
\endfirsthead
\multicolumn{4}{@{}l}{\emph{\tablename~\thetable{} (cont.)}}\\
\toprule
Title & Year & Provider & Reason \\
\midrule
\endhead
\midrule \multicolumn{4}{r@{}}{\emph{continued on next page}}\\
\endfoot
\bottomrule
\endlastfoot
A review of multiple-choice item-writing guidelines for classroom assessment & 2002 & WoS, OA & \emph{cited}: item-writing guidelines \\
A Systematic Assessment of `None of the Above' on Multiple Choice Tests in a First Year Psychology Classroom & 2015 & OA & \emph{cited}: none/all of the above \\
An Examination of the Effect of Multiple-Choice Item Writing Quality on Test Performance: Differential Effects on Subgroups? & 2025 & OA & \emph{cited}: effects of flaws on item statistics \\
Differential Length and Overlap with the Stem in Multiple-Choice Item Options: A Pilot Experiment & 2018 & OA & \emph{cited}: use of cues by examinees \\
Do item-writing flaws reduce examinations psychometric quality? & 2016 & OA & \emph{cited}: effects of flaws on item statistics \\
Do sequentially-presented answer options prevent the use of testwiseness cues on continuing medical education tests? & 2014 & OA & \emph{cited}: use of cues by examinees \\
Evaluation of five guidelines for option development in multiple-choice item-writing & 2009 & WoS, OA & \emph{cited}: use of cues by examinees \\
Examining the effects of testwiseness in conceptual physics evaluations & 2016 & OA & \emph{cited}: use of cues by examinees \\
Item-Writing Practice and Evidence & 2011 & OA & \emph{cited}: item-writing guidelines \\
Multiple-choice test items with ``none of the above'' as an alternative & 2025 & OA & \emph{cited}: none/all of the above \\
The Art and Science of Item Writing: A Review of Established Guidelines for Multiple-Choice Questions & 2025 & OA & \emph{cited}: item-writing guidelines \\
The Effects of Violating Standard Item Writing Principles on Tests and Students: The Consequences of Using Flawed Test Items on Achievement Examinations in Medical Education & 2005 & OA & \emph{cited}: effects of flaws on item statistics \\
The impact of item-writing flaws and item complexity on examination item difficulty and discrimination value & 2016 & OA & \emph{cited}: effects of flaws on item statistics \\
The impact of item-writing flaws on difficulty and discrimination in item response theory & 2026 & OA & \emph{cited}: effects of flaws on item statistics \\
           & 2020 & OA & supports the claim (0.75) \\
              & 2019 & OA & supports the claim (0.40) \\
               & 2019 & OA & supports the claim (0.40) \\
                  & 2006 & OA & supports the claim (0.40) \\
                        & 2023 & OA & supports the claim (0.40) \\
                         ( 8990) & 2015 & OA & supports the claim (0.35) \\
  \textquotedbl{}  \textquotedbl{}          & 2020 & OA & supports the claim (0.55) \\
  Writing Multiple-Choice Questions to Promote Critical Thinking Based on Nursing Process & 2014 & OA & supports the claim (0.50) \\
2012-2021 Yılları Arasındaki Tıpta Uzmanlık Sınavındaki Madde Türlerinin İncelenmesi & 2026 & OA & supports the claim (0.45) \\
\textquotedbl{}                           & 2010 & OA & supports the claim (0.60) \\
\textquotedbl{}                    \textquotedbl{} & 2017 & OA & supports the claim (0.70) \\
A Baseline for Multiple-Choice Testing in the University Classroom & 2021 & OA & supports the claim (0.55) \\
A Comparison of Item Difficulty and Item Discrimination of Multiple-Choice Items Using \textquotedbl{}None of the Above\textquotedbl{} Options. & 1986 & OA & supports the claim (0.80) \\
A COMPARISON OF THE ITEM DIFFICULTY AND ITEM DISCRIMINATION OF MULTIPLE-CHOICE ITEMS USING THE NONE OF THE ABOVE AND ONE CORRECT RESPONSE OPTIONS & 1987 & WoS, OA & supports the claim (0.55) \\
A Comparison of the Quality of Multiple-Choice Questions from CPA Exams and Textbook Test Banks & 1998 & OA & supports the claim (0.70) \\
A few considerations of climbing-snail problem: Fibonacci's error, problem's popularity and Mexican students' performances & 2017 & OA & supports the claim (0.50) \\
A Guide for Writing Single Best Answer Questions to Assess Higher-Order Thinking Skills Based on Learning Outcomes & 2022 & OA & supports the claim (0.60) \\
A Guideline on Developing Effective Multiple Choice Questions and Construction of Single Best Answer format & 2012 & OA & supports the claim (0.70) \\
A Novel Method for Evaluating Examination Item Quality & 2015 & OA & supports the claim (0.85) \\
A practical discussion to avoid common pitfalls when constructing multiple choice questions items & 2010 & OA & supports the claim (0.90) \\
A preventive intervention in the Language Testing and Assessment course: building pre-service teachers' multiple-choice item writing skills & 2026 & OA & supports the claim (0.45) \\
A Primer for Developing Measures of Science Content Knowledge for Small-Scale Research and Instructional Use & 2016 & OA & supports the claim (0.40) \\
A Simple Methodology for Discerning Item Construction Flaws in Health Professions Examinations & 2018 & OA & supports the claim (0.75) \\
A Study of Writing Chinese Language Test Items : Based on The Critical Evaluation of MCQ Items in New HSK, SNULT and FLEX & 2015 & OA & supports the claim (0.45) \\
A Taxonomy of Multiple-Choice Item-Writing Rules & 1989 & OA & supports the claim (0.80) \\
Adapting multiple-choice item-writing guidelines to an industrial context & 2010 & Scopus, WoS, OA & supports the claim (0.70) \\
AI versus human-generated multiple-choice questions for medical education: a cohort study in a high-stakes examination & 2025 & OA & supports the claim (0.40) \\
Alignment of Classroom History Assessments and the 7th National Curriculum in Korea: Assessing Historical Knowledge and Reasoning Skills & 2005 & OA & supports the claim (0.60) \\
An analysis of multiple choice items of Computer Based English Test at SMP AlKhairiyah & 2019 & OA & supports the claim (0.50) \\
An Analysis of the Quality of Multiple-Choice Test Items in the Final Summative Assessment for Grade X Indonesian Language Learning & 2026 & OA & supports the claim (0.45) \\
AN ANALYSIS OF THE QUALITY OF TEACHER -- MADE MULTIPLE CHOICE TESTS USED AS SUMMATIVE ASSESSMENT FOR ENGLISH SUBJECT AT SMP NEGERI 5 SINGARAJA & 2020 & OA & supports the claim (0.50) \\
An Assessment of the Effectiveness of Complex Alternatives in Multiple Choice Achievement Test Items & 1975 & OA & supports the claim (0.90) \\
An Automatic Question Usability Evaluation Toolkit & 2024 & OA & supports the claim (0.40) \\
ANALISIS BUTIR SOAL PILIHAN GANDA MATA PELAJARAN BAHASA INDONESIA SMK NEGERI 3 SINGARAJA & 2014 & OA & supports the claim (0.40) \\
Analisis butir soal ulangan akhir semester mata pelajaran Pendidikan Agama Islam (PAI) kelas VIII semester genap sekolah menengah pertama (SMP) Kota Semarang tahun pelajaran 2016/ 2017 & 2018 & OA & supports the claim (0.35) \\
Analisis Konstruksi Butir Soal Pilihan Ganda Berdasarkan Kaidah Penulisan yang Baik & 2026 & OA & supports the claim (0.45) \\
Analysis of item writing flaws (IWFs) evident in objective formats examination questions in federal college of education (technical) Asaba, Nigeria & 2021 & OA & supports the claim (0.50) \\
Analysis of Item Writing Flaws in a Communications Skills Test in a Ghanaian University & 2021 & OA & supports the claim (0.50) \\
Analysis of medical and dentistry basic sciences examinations: A case study & 2020 & OA & supports the claim (0.35) \\
Analysis of the Orthodontic Exams in Department of Orthodontics of Shahid Beheshti Dental School & 2021 & OA & supports the claim (0.50) \\
ANALYSIS OF THE QUALITY OF SUMMATIVE ARABIC TEST ITEMS FOR GRADE V AT AT-TAQWA ISLAMIC ELEMENTARY SCHOOL BASED ON HALADYNA'S THEORY & 2026 & OA & supports the claim (0.55) \\
Anatomy of a Clerkship Test & 2010 & OA & supports the claim (0.60) \\
Análise da adequação dos itens do Teste de Progresso em medicina & 2022 & OA & supports the claim (0.45) \\
ANÁLISE DE QUESTÕES DE MÚLTIPLA ESCOLHA E PROPOSTA DE ESTRATÉGIAS PARA A MELHORIA DA ELABORAÇÃO DE AVALIAÇÕES COGNITIVAS DO CURSO DE MEDICINA DA UNIVERSIDADE MUNICIPAL DE SÃO CAETANO DO SUL - CAMPUS SÃO PAULO & 2025 & OA & supports the claim (0.45) \\
Análisis de validez de la prueba Planea de EMS 2017, área matemática & 2017 & OA & supports the claim (0.60) \\
Applying the utility index to review single best answer questions in medical education assessment & 2020 & OA & supports the claim (0.40) \\
Aspetti di criticità delle prove contenute nel volume ``il test di italiano per stranieri'' & 2015 & OA & supports the claim (0.60) \\
ASSESSING INTER-RATER AGREEMENT ABOUT ITEM-WRITING FLAWS IN MULTIPLE-CHOICE QUESTIONS OF CLINICAL ANATOMY & 2013 & WoS, OA & supports the claim (0.55) \\
Assessing multiple-choice questions based on language precision and best practices to promote equity in the Dental Physiology course & 2023 & OA & supports the claim (0.80) \\
Assessing the Knowledge of Teachers in Objective Test Construction Procedure in the Teacher Education Programs: A Literature Review & 2023 & OA & supports the claim (0.45) \\
Assessing the Test-Wiseness of Health Science Students: Development and Validation of an Instrument & 1985 & OA & supports the claim (0.75) \\
Assessing The Use of High Quality Multiple Choice Exam Questions in Undergraduate Nursing Education: Are Educators Making the Grade? & 2011 & OA & supports the claim (0.70) \\
Assessment by modified essay type questions and multiple-choice questions in medical undergraduate students: Which are better in addressing higher cognitive skills? & 2024 & OA & supports the claim (0.40) \\
Assessment of distractor efficiency of MCQS in item analysis. & 2022 & OA & supports the claim (0.55) \\
Assessment of financial accounting questions in the May/June West African Senior School Certificate Examination from 2014 to 2024 & 2026 & OA & supports the claim (0.70) \\
Assessment of Item-Writing Flaws in Multiple-Choice Questions & 2013 & OA & supports the claim (0.75) \\
Assessment of Item-Writing Flaws in Multiple-Choice Questions & 2013 & OA & supports the claim (0.40) \\
Assessment of multiple-choice questions by item analysis for medical students' examinations & 2022 & OA & supports the claim (0.55) \\
Assessment of multiple-choice questions in selected test banks accompanying text books used in nursing education & 2001 & Scopus, WoS, OA & supports the claim (0.80) \\
AUF DER SUCHE NACH EINEM OPTIMALEN KONZIPIERUNGSVERFÄHREN DER TESTAUFGABEN ZUR ÜBERPRÜFUNG DER SCHÜLERLEISTUNGEN & 1991 & OA & supports the claim (0.40) \\
Back to basics: reflective take on role of MCQs in undergraduate Malaysian dental professional qualifying exams & 2023 & OA & supports the claim (0.45) \\
BenchMarker: An Education-Inspired Toolkit for Highlighting Flaws in Multiple-Choice Benchmarks & 2026 & OA & supports the claim (0.40) \\
Best Practices for Writing Test Items & 2012 & WoS, OA & supports the claim (0.50) \\
Best Practices on Examination Construction, Administration, and Feedback & 2018 & OA & supports the claim (0.65) \\
Best Practices Related to Examination Item Construction and Post-hoc Review & 2019 & OA & supports the claim (0.75) \\
Bones pràctiques en l'ús de proves d'alternativa múltiple & 2019 & OA & supports the claim (0.60) \\
Buenas prácticas en el uso de pruebas de alternativa múltiple & 2019 & OA & supports the claim (0.65) \\
Choosing medical assessments: Does the multiple-choice question make the grade? & 2018 & OA & supports the claim (0.75) \\
CME activities of medical journals:: Quality of multiple-choice questions as evaluation tool.: Using the example of the German medical journals Deutsches Arzteblatt, Deutsche Medizinische Wochenschrift, and Der Internist & 2007 & WoS & supports the claim (0.55) \\
Cognitive-Level and Item-Quality Analysis of Multiple-Choice Items inan Introductory Oil, Gas, and Renewable Energy Course & 2026 & OA & supports the claim (0.40) \\
Collaborative Examination Item Review Process in a Team-Taught, Self-Care Sequence & 2015 & OA & supports the claim (0.35) \\
Common Item Violations in Multiple Choice Questions in Bangladeshi Recruitment Tests & 2022 & OA & supports the claim (0.50) \\
Common item-writing flaws and psychometric properties of multiple-choice test in technical schools in Ghana & 2026 & OA & supports the claim (0.75) \\
Como elaborar questões de múltipla escolha de boa qualidade & 2019 & OA & supports the claim (0.70) \\
Comparison of emergency medical residency pre-board and board exams among universities of medical sciences in Iran & 2020 & OA & supports the claim (0.40) \\
Competence and Common Errors in Multiple-Choice Test Construction Among Junior High School English Teachers: Foundation for a Proposed Intervention Program & 2026 & OA & supports the claim (0.45) \\
Computer-assisted item and test pre-analysis: a new direction in qualitative methods & 1987 & OA & supports the claim (0.45) \\
Construct-irrelevant variance and flawed test questions: Do multiple-choice item-writing principles make any difference? & 2002 & WoS, OA & supports the claim (0.85) \\
Constructing Multiple-Choice Items to Measure Higher-Order Thinking & 2017 & OA & supports the claim (0.45) \\
Constructing Test Items: Multiple-Choice, Constructed-Response, Performance and Other Formats & 2006 & OA & supports the claim (0.40) \\
Constructing valid classroom reading test tasks (multiple choice questions) for university students (notes on a workshop) & 2015 & OA & supports the claim (0.45) \\
Construction and Writing Flaws of the Multiple-Choice Questions in the Published Test Banks of Obstetrics and Gynecology: Adoption, Caution, or Mitigation? & 2022 & OA & supports the claim (0.75) \\
Construction of Multiple Choice Questions Before and After An Educational Intervention & 2017 & OA & supports the claim (0.45) \\
Crafting the Question: Get the Most Out of Your Student Response System & 2016 & OA & supports the claim (0.50) \\
Creating High Quality Single Best Answer Questions: A Guide for Medical Educators & 2024 & OA & supports the claim (0.45) \\
Creating multiple-choice items for testing student learning & 2022 & OA & supports the claim (0.85) \\
Cross-cultural study of test-wiseness & 2010 & OA & supports the claim (0.60) \\
Crowdsourcing the Evaluation of Multiple-Choice Questions Using Item-Writing Flaws and Bloom's Taxonomy & 2023 & OA & supports the claim (0.45) \\
Cueing effects of item writing flaws in multiple-choice questions & 2013 & OA & supports the claim (0.95) \\
Cuing effect of \textquotedbl{}all of the above\textquotedbl{} on the reliability and validity of multiple-choice test items & 1998 & WoS, OA & supports the claim (0.85) \\
Damage control orthopedic surgery (DOC): Is there an Influence on outcome? & 2011 & OA & supports the claim (0.70) \\
Detecting Flaws in this Old Test & 2023 & OA & supports the claim (0.90) \\
Determining the alignment of assessment items with curriculum goals through document analysis by addressing identified item flaws & 2025 & OA & supports the claim (0.75) \\
Developing an Automated Evaluation Tool for Multiple-Choice Questions & 2023 & Scopus & supports the claim (0.50) \\
Developing and validating multiple-choice test items, third edition & 2012 & Scopus & supports the claim (0.95) \\
Developing multiple - choice questions in mathematics & 2009 & OA & supports the claim (0.50) \\
Development and Assessment of a Multiple-Choice Quiz for a Bending Stress Quiz & 2025 & OA & supports the claim (0.35) \\
Development and content validation of the NEIA tool: An instrument to evaluate NCLEX-style exam items & 2026 & OA & supports the claim (0.35) \\
Development and use of a multiple-choice item writing flaws evaluation instrument in the context of general chemistry & 2019 & WoS & supports the claim (0.70) \\
Development of multiple choice questions' instructional plan for the nursing educators & 2019 & OA & supports the claim (0.50) \\
Differences in Multiple-Choice Questions of Opposite Stem Orientations Based on a Novel Item Quality Measure & 2022 & OA & supports the claim (0.45) \\
Differential Effects of the Discrete Options Multiple Choice Test Format on Test Takers Assesment and Scores & 2010 & OA & supports the claim (0.45) \\
Directrices para la construcción de ítems de elección múltiple & 2004 & OA & supports the claim (0.55) \\
Distractor similarity and item-stem structure: Effects on item difficulty & 2007 & WoS, OA & supports the claim (0.70) \\
Document-Based Assessment of Item-Writing Skills: Implications for Policy Formulation in Teacher Education & 2026 & OA & supports the claim (0.65) \\
Does central vetting improve construct quality of one-best-answer items in medical school: An audit & 2017 & OA & supports the claim (0.65) \\
Does Educator Training or Experience Affect the Quality of Multiple-Choice Questions? & 2015 & OA & supports the claim (0.35) \\
Does the Response Options Placement Provide Clues to the Correct Answers in Multiple-choice Tests? A Systematic Review & 2022 & OA & supports the claim (0.60) \\
E-workshop on writing multiple-choice questions: Online experience during the period of COVID-19 crisis & 2020 & OA & supports the claim (0.45) \\
Educator's blueprint: A howto guide for developing highquality multiplechoice questions & 2023 & OA & supports the claim (0.60) \\
Effect of changing multiple choice questions from ``all of the above'' to ``select all that apply'' & 2024 & OA & supports the claim (0.80) \\
Effect of Faculty Training Programs on Improving Quality of Residency Exams in 2013-2014 & 2015 & OA & supports the claim (0.65) \\
EFFECT OF ITEM WRITING FLAWS IN MULTIPLE CHOICE QUESTIONS ON STUDENTS' ACHIEVEMENT GROUPS AND RELIABILITIES OF TESTS IN AZAD JAMMU KASHMIR MEDICAL COLLEGE & 2024 & OA & supports the claim (0.85) \\
Effect of Technically Worded Options on Multiple-Choice Test Performance. & 1980 & OA & supports the claim (0.85) \\
EFFECTIVENESS OF A FACULTY DEVELOPMENT PROGRAM ON THE QUALITY OF A-TYPE MULTIPLE CHOICE QUESTIONS IN UNDERGRADUATE MEDICAL AND DENTAL EDUCATION: A QUASI-EXPERIMENTAL STUDY & 2026 & OA & supports the claim (0.55) \\
Effectiveness of longitudinal faculty development programs on MCQs items writing skills: A follow-up study & 2017 & OA & supports the claim (0.60) \\
EFFECTIVENESS OF OPTIONS IN MULTIPLE-CHOICE ITEMS: THE CASE OF ``NONE OF THE ABOVE'' & 2015 & OA & supports the claim (0.70) \\
Effects of a long term faculty development program on improvement in quality of MCQs: an impact evaluation study & 2025 & OA & supports the claim (0.50) \\
Effects of Mayfield's Four Questions (M4Q) on nursing students' self-efficacy and multiple-choice test scores & 2010 & OA & supports the claim (0.45) \\
Effects of multiple-choice item-writing guideline utilization on item and student performance & 2013 & OA & supports the claim (0.75) \\
Effects of Noncontent Variables on Mathematics Test Performance & 1987 & OA & supports the claim (0.75) \\
Effects of Question Formats on Student and Item Performance & 2013 & WoS, OA & supports the claim (0.85) \\
Effects of Reducing the Cognitive Load of Mathematics Test Items on Student Performance & 2015 & OA & supports the claim (0.40) \\
Effects of Structural Flaws on the Psychometric Properties of Multiple-Choice Questions & 2018 & OA & supports the claim (0.80) \\
EFFECTS OF TESTWISENESS STRATEGY ON THE ACADEMIC PERFORMANCE OF SENIOR SECONDARY SCHOOL STUDENTS IN ECONOMICS IN OYO STATE, NIGERIA & 2022 & OA & supports the claim (0.40) \\
Efficacy of Faculty Development Training Workshops (FDTWs) on Writing High-Quality Multiple-Choice Questions at Northern Border University (NBU) in the Kingdom of Saudi Arabia (KSA) & 2024 & OA & supports the claim (0.50) \\
Ein Plädoyer zur Qualitätssicherung schriftlicher Prüfungen im Psychologiestudium & 2021 & OA & supports the claim (0.60) \\
Elimination of Distractors: A Construct-irrelevant Strategy? An Investigation of Examinees' Response Decision Processes in an EFL Multiple-choice Reading Test & 2018 & OA & supports the claim (0.45) \\
ENHANCING MULTIPLE-CHOICE TEST CONSTRUCTION: AN EXPERIENTIAL LEARNING FOR BIOLOGY TEACHERS & 2016 & OA & supports the claim (0.35) \\
Errors in the construction of multi-choice questions: An analysis & 2018 & OA & supports the claim (0.70) \\
Establishing the validity and reliability of the fairness of items tool & 2014 & OA & supports the claim (0.70) \\
Evaluating single-best-answer questions in dental materials education: a retrospective item analysis from an online undergraduate dentistry question bank & 2026 & OA & supports the claim (0.55) \\
Evaluating the Effectiveness of a Faculty Development Program in Enhancing the Quality of Multiple-Choice Questions: A Quantitative Study & 2026 & OA & supports the claim (0.60) \\
Evaluating the Impact of Technical Flaws in Multiple-Choice Questions on Student Performance in Physiology & 2025 & OA & supports the claim (0.85) \\
Evaluating the Psychometric Qualities of Higher Education Instructor-constructed Multiple-Choice Tests: Impact on Test Quality and Consequences for Student Grades & 2013 & OA & supports the claim (0.65) \\
Evaluating the Quality of Examination Items From the Pathophysiology, Drug Action, and Therapeutics Course Series & 2024 & OA & supports the claim (0.40) \\
Evaluating the Quality of Higher Education Instructor-Constructed Multiple-Choice Tests: Impact on Student Grades & 2017 & OA & supports the claim (0.50) \\
Evaluating the Quality of Multiple Choice Question in Paediatric Dentistry Postgraduate Examinations & 2019 & OA & supports the claim (0.65) \\
Evaluating the Quality of Multiple-Choice Assessment: Evidence from Bhutan & 2026 & OA & supports the claim (0.50) \\
Evaluating the quality of multiple-choice questions in a NAPLEX preparation book & 2020 & OA & supports the claim (0.70) \\
Evaluating the quality of multiplechoice question pilot database: A global educatorcreated tool for conceptbased pharmacology learning & 2024 & OA & supports the claim (0.40) \\
Evaluation and Improvement of Multiple-choice Questions in the Doctor of Pharmacy Program & 2020 & OA & supports the claim (0.65) \\
Evaluation of cognitive levels and item writing flaws in medical pharmacology internal assessment examinations & 2017 & OA & supports the claim (0.65) \\
Evaluation of MCQs from MOOCs for common item writing flaws & 2018 & OA & supports the claim (0.70) \\
Evaluation of medical assays by multiple-choice questions: Example of 10-years' experience of a North African Country & 2022 & OA & supports the claim (0.40) \\
Evaluation of Modified Essay Questions (MEQ) and Multiple Choice Questions (MCQ) as a tool for Assessing the Cognitive Skills of Undergraduate Medical Students & 2011 & WoS & supports the claim (0.40) \\
Evaluation of Multiple Choice and Short Essay Question items in Basic Medical Sciences & 2014 & WoS, OA & supports the claim (0.60) \\
Evaluation of Multiple Choice Questions of Oral and Maxillofacial Medicine Courses 1, 2, and 3 in the First Semester of Academic Year 2014-2015 & 2018 & OA & supports the claim (0.80) \\
Evaluation of multiple-choice and short essay questions in pharmacology education & 2023 & OA & supports the claim (0.60) \\
Evaluation of Pharmacology MCQ Papers of MBBS Professional Examinations & 2022 & OA & supports the claim (0.55) \\
Exam theory for dummies & 2020 & OA & supports the claim (0.50) \\
Examining Gender Bias in Multiple Choice Item Formats Violating Item-Writing Guidelines* & 2019 & OA & supports the claim (0.60) \\
Examining Multiple-choice Items in High School English Tests---Linking to Assessment Literacy & 2015 & OA & supports the claim (0.50) \\
Examining Social Studies Course Questions in Terms of Question Type and Test Item Writing Techniques & 2023 & OA & supports the claim (0.55) \\
Examining the impact of specific types of item-writing flaws on student performance and psychometric properties of the multiple choice question & 2018 & OA & supports the claim (0.85) \\
Exploring student perspectives and factors that may influence performance in multiple-choice physiology assessments & 2025 & OA & supports the claim (0.85) \\
Exploring the psychometric properties of the Testwiseness Rating Scale using a measurement modeling approach & 2026 & OA & supports the claim (0.45) \\
Factor Structure of the Gibb Experimental Test of Testwiseness & 1990 & OA & supports the claim (0.60) \\
Faculty development programs improve the quality of Multiple Choice Questions items' writing & 2015 & OA & supports the claim (0.60) \\
Fit for purpose? Evaluating multiple-choice question quality in E-learning for emergency primary healthcare & 2024 & OA & supports the claim (0.45) \\
Flawed multiple-choice questions put on the scale: What is their impact on students' achievement in a final undergraduate surgical examination? & 2016 & OA & supports the claim (0.85) \\
Flaws of Multiple Choice Questions in Teacher-Constructed Nursing Examinations: A Pilot Descriptive Study & 2017 & OA & supports the claim (0.70) \\
GOOD MULTIPLE CHOICE CRITERIA: COMPARISON BETWEEN ENGLISH TEACHERS' UNDERSTANDING AND TAXONOMY OF MULTIPLE-CHOICE ITEM-WRITING GUIDELINES & 2017 & OA & supports the claim (0.45) \\
Guess for Success? Application of a Mixture Model to Test-Wiseness on Multiple-Choice Exams & 2023 & OA & supports the claim (0.60) \\
Guide To Developing High-Quality, Reliable, and Valid Multiple-Choice Assessments & 2014 & OA & supports the claim (0.55) \\
Guidelines based on validity criteria for the development of multiple choice items & 2015 & OA & supports the claim (0.55) \\
Guidelines for Creating Written Clinical Reasoning Exams: Insight from a Delphi Study & 2018 & OA & supports the claim (0.70) \\
Guidelines for the construction of multiple choice questions tests & 2006 & OA & supports the claim (0.55) \\
Guidelines for the design of digital closed questions for assessment and learning in higher education & 2007 & OA & supports the claim (0.65) \\
Guidelines for Writing Multiple-Choice Questions in Radiology Courses & 1986 & OA & supports the claim (0.90) \\
Guidelines on Writing Multiple Choice Questions: A Well-Received and Effective Faculty Development Intervention & 2020 & OA & supports the claim (0.55) \\
Guiding Principles to Writing Accessible Multiple-Choice Questions (MCQs): A practical evidence-based guide for educators & 2026 & OA & supports the claim (0.60) \\
Gute Fragen für gute Lehre. Eine Handreichung für Lehrende zur Erstellung von Fragen im Antwortwahlverfahren. & 2020 & OA & supports the claim (0.80) \\
How (not) to write an introductory programming exam & 2015 & OA & supports the claim (0.60) \\
How Do I Create Good Constructed-Response Items? & 2015 & OA & supports the claim (0.35) \\
How Do You Know if They're Getting It? Writing Assessment Items that Reveal Student Understanding. & 2009 & OA & supports the claim (0.50) \\
How-to-guide for writing multiple choice questions for the pharmacy instructor & 2017 & WoS & supports the claim (0.60) \\
Identification of technical item flaws leads to improvement of the quality of single best Multiple Choice Questions & 2013 & OA & supports the claim (0.60) \\
Identifying item flaws to improve the quality of core mathematics Multiple-Choice Items & 2024 & OA & supports the claim (0.65) \\
Identifying the presence of context and item-writing flaws in practice items: The case of Turkish mathematics textbooks & 2022 & OA & supports the claim (0.50) \\
Impact Measurement on Medical Faculty for Adhering to Appropriate Guidelines in Framing Effective Multiple-Choice Questions for Item Analysis & 2020 & OA & supports the claim (0.50) \\
Impact of Flawed Multiple-Choice Questions on Tests and Student Achievements in Medical Education & 2019 & OA & supports the claim (0.85) \\
Improving Nursing Faculty Comprehension of Evaluation Methods & 2018 & OA & supports the claim (0.40) \\
Improving the Quality of Multiple-Choice Questions (MCQs) in Teacher Written Tests: Experimenting with Malaysian Science Teachers & 2025 & OA & supports the claim (0.50) \\
INOValiar: lista consolidada de critérios de qualidade de questões de escolha múltipla na avaliação de aprendizagens & 2025 & OA & supports the claim (0.75) \\
Inovasi Kaidah Penulisan Soal Pilihan Ganda dalam Evaluasi Pembelajaran: Tantangan Literasi dan Penalaran Peserta Didik & 2026 & OA & supports the claim (0.35) \\
Investigating possible causes of bias in a progress test translation: an one-edged sword & 2019 & OA & supports the claim (0.45) \\
Investigating the Impact of Item Writing Flaws on Student Performance in Physiology Assessments at Two Institutions & 2026 & OA & supports the claim (0.85) \\
Investigating the Influence of Proficiency and Gender on the Use of Selected Test-Wiseness Strategies in Higher Education & 2008 & OA & supports the claim (0.70) \\
Investigating the potential influence of established multiple-choice test-taking cues on item response in a pharmacotherapy board certification examination preparatory manual: A pilot study & 2006 & WoS, OA & supports the claim (0.80) \\
Is the formal quality of multiple choice questions used for continuing medical education influenced by guidelines for authors provided by the publishing houses? & 2017 & WoS & supports the claim (0.50) \\
ISSUES AND ADVANCEMENTS IN ASSESSMENT: A SCOPING REVIEW OF MULTIPLE-CHOICE QUESTIONS IN MEDICAL EDUCATION & 2026 & OA & supports the claim (0.45) \\
Item analysis and optimizing multiple-choice questions for a viable question bank in ophthalmology: A cross-sectional study & 2021 & OA & supports the claim (0.50) \\
Item Analysis of Multiple Choice Questions in Pharmacology & 2020 & OA & supports the claim (0.60) \\
Item Analysis: The impact of distractor efficiency on the discrimination power of multiple choice items & 2019 & OA & supports the claim (0.45) \\
Item Development & 2021 & OA & supports the claim (0.35) \\
Item Writing for Nurse Educators & 2023 & OA & supports the claim (0.70) \\
Item-writing flaws and answer-position bias in large language model-generated medical multiple-choice questions: a benchmark study of 15 models with an open corpus and a rule-based detector & 2026 & OA & supports the claim (0.40) \\
Item-Writing Guidelines on Response Option Placement: A Systematic Review & 2024 & OA & supports the claim (0.60) \\
Item-writing rules: Collective wisdom & 2005 & WoS & supports the claim (0.70) \\
Knowledge of dental faculty in gulf cooperation council states of multiple-choice questions' item writing flaws & 2020 & OA & supports the claim (0.55) \\
Knowledge, application and how about competence? Qualitative assessment of multiple-choice questions for dental students & 2020 & OA & supports the claim (0.50) \\
Mastering the Art of Framing Multiple-choice Questions and Improving Quality of the Same through Medical Faculty Training & 2026 & OA & supports the claim (0.55) \\
MCQ MULTIPLE CHOICE QUESTIONS EXAMINATION WITHOUT ERRORS IN PHARMACOLOGY & 2021 & OA & supports the claim (0.75) \\
Millî Eğitim Bakanlığı Tarafından Yayınlanan Test Maddelerinin Madde Yazım İlkelerine Uygunluğunun İncelenmesi & 2022 & OA & supports the claim (0.65) \\
Mnemonic strategy for multiple choice question formulation- quick guide to medical educators & 2022 & OA & supports the claim (0.80) \\
More Multiple-Choice Item Writing Do's and Don'ts & 1994 & OA & supports the claim (0.55) \\
More Multiple-Choice Item Writing Do's and Don'ts. ERIC/AE Digest. & 1995 & OA & supports the claim (0.55) \\
MULTI-CHOICE QUESTIONS AND THEIR PROBLEMS WHEN USED FOR ASSESSMENT OF AIRCRAFT ENGINEERS EDUCATION & 2017 & OA & supports the claim (0.60) \\
Multiple Choice Question Metamorphosis---Enhancing Question Quality through Faculty Training: A Pragmatic Study & 2026 & Scopus & supports the claim (0.55) \\
Multiple Choice Questions in Medical Education: How to Construct High Quality Questions & 2020 & OA & supports the claim (0.65) \\
MULTIPLE CHOICE QUESTIONS; DEVELOPED BY THE FACULTY OF A PUBLIC SECTOR MEDICAL COLLEGE & 2017 & OA & supports the claim (0.50) \\
Multiple Choice Test-Taking Strategies, Test Anxiety, and EFL Students' Achievement & 2023 & OA & supports the claim (0.40) \\
Multiple-choice questions : linguistic investigation of difficulty for first-language and second-language students & 2010 & OA & supports the claim (0.50) \\
Multiple-choice questions : teaching guide for higher and professional education & 2019 & OA & supports the claim (0.45) \\
Multiple-choice test development competencies of junior high school mathematics teachers in Ghana: a triangulation methodology & 2025 & OA & supports the claim (0.40) \\
Multiple-choice tests: a tool in assessing knowledge & 2010 & OA & supports the claim (0.40) \\
Multiple-Choice-Aufgaben : Teaching Guide for Higher \& Professional Education & 2018 & OA & supports the claim (0.60) \\
Negatively-Worded Multiple Choice Questions: An Avoidable Threat to Validity & 2017 & OA & supports the claim (0.60) \\
New Guidelines for Developing Multiple-Choice Items & 2006 & OA & supports the claim (0.60) \\
New test-taking patterns and their effect on language test validity & 2019 & OA & supports the claim (0.60) \\
Nonfunctional distractor analysis: An indicator for quality of Multiple choice questions & 2020 & OA & supports the claim (0.60) \\
NYSE Multiple-Choice Item Writing Manual & 1999 & OA & supports the claim (0.45) \\
Námsmatsaðferðir - Fjölvalspróf & 2019 & OA & supports the claim (0.55) \\
Optimising Medical Education: The Power and Potential of One Best Answer (OBA) Questions & 2024 & OA & supports the claim (0.50) \\
Option Guessability Assessment and Item-Writing Quality Control;An Empirical Study on China's NMEE MCQ Options & 2012 & OA & supports the claim (0.75) \\
Passage Dependence of Multiple-Choice Items in Reading Comprehension Component of English Academic Test of Guangdong Province & 2019 & OA & supports the claim (0.35) \\
PASSAGE DEPENDENCE OF THE NEW SAT® READING COMPREHENSION QUESTIONS & 1993 & OA & supports the claim (0.70) \\
Position of Correct Option and Distractors Impacts Responses to Multiple-Choice Items: Evidence From a National Test & 2022 & OA & supports the claim (0.55) \\
Post exam item analysis: Implication for intervention & 2019 & OA & supports the claim (0.65) \\
Postexamination item analysis of undergraduate pediatric multiple-choice questions exam: implications for developing a validated question Bank & 2024 & OA & supports the claim (0.45) \\
Postgraduate Clinical Residency: The Impact of Multiple-Choice Question Quality on Exam Success Rates & 2025 & OA & supports the claim (0.70) \\
Pre and Post examination evaluation of MCQs and SEQs Items in Undergraduate MBBS mock examination. & 2020 & OA & supports the claim (0.45) \\
Prediction of item psychometric indices from item characteristics automatically extracted from the stem and option text & 2011 & WoS & supports the claim (0.70) \\
Prevalence of Flawed Multiple-Choice Questions in Continuing Medical Education Activities of Major Radiology Journals & 2015 & WoS & supports the claim (0.50) \\
Primary school teachers' competences in constructing quality multiple-choice questions in Tanzania & 2025 & OA & supports the claim (0.70) \\
Problems Faced by English Teacher Candidates in Developing Test Kits for Assessing Students' Learning & 2023 & OA & supports the claim (0.50) \\
Protocolo De Intervención Para La Construcción Y Validación De Ítems De Pruebas Funcionales Escritas En Un Concurso De Méritos & 2021 & OA & supports the claim (0.45) \\
Psychometric item/question analysis of multiple-choice questions in fixed prosthodontics exam & 2025 & OA & supports the claim (0.60) \\
Psychometrics and test-enhanced learning in a patient-centered learning curriculum & 2014 & OA & supports the claim (0.60) \\
Qualidade dos itens de múltipla escolha utilizados em um Teste de Progresso & 2022 & OA & supports the claim (0.45) \\
Quality assurance of item writing: During the introduction of multiple choice questions in medicine for high stakes examinations & 2008 & OA & supports the claim (0.70) \\
Quality Multiple-Choice Test Questions: Item-Writing Guidelines and an Analysis of Auditing Testbanks & 1997 & OA & supports the claim (0.60) \\
Quality of design and analysis of pre-internship exams for medical students of Hamadan University of Medical Sciences in the years 2018 to 2020 & 2022 & OA & supports the claim (0.50) \\
Question Design for Classroom Voting in Undergraduate Mathematics: A Systematic Review of the Literature & 2025 & OA & supports the claim (0.80) \\
Question vetting: theory and practice & 2012 & OA & supports the claim (0.40) \\
Questions for Classroom Response Systems and Teaching Instrumental Element Analysis & 2021 & OA & supports the claim (0.60) \\
Racial Politics and the Invisibility of Asian Americans in Higher Education. & 1995 & OA & supports the claim (0.80) \\
Reconstructing a Junior High School English Question Bank through Qualitative Item Analysis and Higher-Order Thinking Skills Integration & 2026 & OA & supports the claim (0.50) \\
Refine Test Items for Accurate Measurement: Six Valuable Tips & 2015 & OA & supports the claim (0.80) \\
Reflexiones sobre adivinar en preguntas de opción múltiple y cómo afecta el resultado del examen & 2015 & OA & supports the claim (0.40) \\
Relations of the Number of Functioning Distractors With the Item Difficulty Index and the Item Discrimination Power in the Multiple Choice Questions & 2023 & OA & supports the claim (0.40) \\
Relationship between Secondary School Students'Level of Testwiseness and their Performance in Mathematics in Bauchi Metropolis of Bauchi State, Nigeria & 2014 & OA & supports the claim (0.40) \\
Relationship between types of distractor and difficulty of multiple-choice vocabulary tests in sentential context & 2013 & OA & supports the claim (0.45) \\
Relationship of text length of multiple-choice questions on item psychometric properties -- A retrospective study & 2020 & OA & supports the claim (0.50) \\
Response Biases in Multiple-Choice Test Item Files & 1982 & OA & supports the claim (0.90) \\
Role of Effective Feed Back in Multiple Choice Questions (MCQs) Designing for Faculty Development & 2010 & OA & supports the claim (0.65) \\
Role of Faculty Development Workshop for Improving MCQS Quality in Basic Medical Sciences & 2021 & OA & supports the claim (0.70) \\
Screening for technical flaws in multiple-choice items. A generalizability study. & 2019 & OA & supports the claim (0.60) \\
Seçme ve Katkı gerektiren Maddelerin Yazımı ve Düzenlenmesi için Kontrol Listeleri & 2023 & OA & supports the claim (0.70) \\
SHORTTERM INSTRUCTION, TESTWISENESS, AND THE SCHOLASTIC APTITUDE TEST: A LITERATURE REVIEW WITH RESEARCH RECOMMENDATIONS & 1978 & OA & supports the claim (0.45) \\
Simplified Guidelines for Multiple-Choice Question Writing to Increase Faculty Compliance and Ensure Valid Student Results & 2019 & OA & supports the claim (0.85) \\
Stability of the Gibb (1964) Experimental Test of Testwiseness & 1988 & OA & supports the claim (0.50) \\
Strategies Used to Improve Quality of Multiple-Choice Questions: A Systematic Review & 2023 & OA & supports the claim (0.60) \\
Structured Process For Writing, Revising, And Assessing Multiple Choice Quizzes & 2020 & OA & supports the claim (0.35) \\
Subjective quality of multiple choice questions used in undergraduate courses in orthopedics and other specialties & 2020 & OA & supports the claim (0.55) \\
Summative Assessment of Cognitive Skills: an Experience Involving Good Practices for Writing Multiple Choice Tests and Exam Composition & 2018 & OA & supports the claim (0.50) \\
Survival of the Fittest: The Role of Linguistic Modification in Nursing Education & 2015 & OA & supports the claim (0.55) \\
Teacher-Made Tests in History & 1979 & OA & supports the claim (0.45) \\
Technical flaws in multiple-choice questions in the access exam to medical specialties (``examen MIR'') in Spain (2009--2013) & 2016 & OA & supports the claim (0.60) \\
Test Construction Principles Violated in Selected Summative Tests in English & 2022 & OA & supports the claim (0.70) \\
Test Question Editing Workshop - Improve Your Test Item Writing Skills Through Editing & 2015 & OA & supports the claim (0.50) \\
Test Wiseness: An Examination of Cue-Using Strategies & 1984 & OA & supports the claim (0.90) \\
Test-Taking Skills and Their Impact on Accessibility for All Students & 2011 & OA & supports the claim (0.40) \\
Test-Taking Skills for Multiple-Choice Formats: Implications for School Psychologists. & 2001 & OA & supports the claim (0.50) \\
Test-Taking Strategies of Arab EFL Learners on Multiple Choice Tests & 2012 & OA & supports the claim (0.45) \\
Test-wiseness : its effect on the supply items of the British Columbia provincial examinations for grade twelve students & 2010 & OA & supports the claim (0.45) \\
Test-wiseness Strategies in PBTs and IBTs: The Case of EFL Test Takers, Who Benefits More? & 2014 & OA & supports the claim (0.55) \\
Test-Wiseness Training: An Investigation of the Impact of Test-Wiseness in an Employment Setting & 2005 & OA & supports the claim (0.65) \\
Test-wiseness: ein unterschätztes Konstrukt? & 2018 & OA & supports the claim (0.70) \\
The ``None of the Above'' Option in Multiple-Choice Testing: An Experimental Study & 2013 & OA & supports the claim (0.85) \\
The Analysis Of The Quality Of Teacher-Made Multiple-Choice Tests Used As Summative Assessment For English Subject & 2020 & OA & supports the claim (0.55) \\
The assessment of quantitative problem-solving skills with \textquotedbl{}none of the above\textquotedbl{} - items (NOTA items) & 2001 & WoS, OA & supports the claim (0.50) \\
The Cognitive Cost of Poor Item Design: How Item-Writing Flaws Shape Response Time and Difficulty Across Assessment Domains & 2026 & OA & supports the claim (0.70) \\
The College Instructor's Guide to Writing Test Items : Measuring Student Learning, by Michael C. Rodriguez and Anthony D. Albano, Routledge, New York: NY, 2017, 173 pp., \$123.90 (Hardcover), ISBN 978-1-138-88653-7 & 2018 & OA & supports the claim (0.70) \\
The Criteria and Analysis of Multiple-Choice Questions in Undergraduate Dental Examinations & 2018 & OA & supports the claim (0.80) \\
The Effect of a One-Day Workshop on the Quality of Framing Multiple Choice Questions in Physiology in a Medical College in India & 2023 & OA & supports the claim (0.55) \\
The Effect of Item Stem and Response Option Length on the Item Analysis Outcomes of a Career and Technical Education Multiple Choice Assessment & 2018 & OA & supports the claim (0.55) \\
The effect of MCQ vetting on students' examination performance & 2014 & OA & supports the claim (0.60) \\
THE EFFECT OF TEST-WISENESS ON SELF-EFFICACY AND MATHEMATIC PERFORMANCE OF MIDDLE SCHOOL STUDENTS WITH LEARNING DISABILITIES & 2011 & OA & supports the claim (0.45) \\
The Effect of Text Length in Crowdsourced Multiple Choice Questions & 2016 & OA & supports the claim (0.50) \\
THE EFFECT OF THE NUMBER OF DISTRACTORS AND THE \textquotedbl{}NONE OF THE ABOVE\textquotedbl{} - \textquotedbl{}ALL OF THE ABOVE\textquotedbl{} OPTIONS IN MULTIPLE CHOICE QUESTIONS & 2021 & OA & supports the claim (0.50) \\
The effect of using \textquotedbl{}none of the above\textquotedbl{} option on psychometric properties of multiple choice questions & 2021 & OA & supports the claim (0.60) \\
The effect of violating an established multiple-choice item writing principle on test difficulty and reliability & 1978 & OA & supports the claim (0.50) \\
The effects of stem completeness and stem orientation on multiple-choice item difficulty and discrimination & 2010 & OA & supports the claim (0.55) \\
The Effects of Test-Taking Practice on Pupil Test Performance & 1968 & OA & supports the claim (0.50) \\
The Effects of Test-Taking Skills and Attitudes on Validity of Standardized Achievement Test Scores of Inner-City Children & 1977 & OA & supports the claim (0.55) \\
The Effects of Test-Wiseness in Medical Education & 1981 & OA & supports the claim (0.70) \\
The Efficacy of Plausible Distracters in the Development of Multiple-Choice Questions & 2026 & OA & supports the claim (0.45) \\
The Extent of Violation of Multiple-Choice Item Writing Rules in the General Secondary Exam in Jordan During Covid-19 & 2024 & OA & supports the claim (0.50) \\
The frequency of item writing flaws in multiple-choice questions used in high stakes nursing assessments & 2006 & WoS, OA & supports the claim (0.60) \\
The future of online testing and assessment: question quality in MOOCs & 2018 & OA & supports the claim (0.60) \\
The Generation and Use of Medical MCQs: A Narrative Review & 2025 & OA & supports the claim (0.70) \\
The impact of a faculty development seminar on the quality of multiple-choice questions & 2016 & WoS & supports the claim (0.40) \\
The Impact of Homogeneity of Answer Choices on Item Difficulty and Discrimination & 2018 & OA & supports the claim (0.80) \\
The Impact of item flaws, testing at low cognitive level, and low distractor functioning on multiple-choice question quality & 2015 & OA & supports the claim (0.80) \\
The Impact Of Item-Writing Flaws On Student Test Scores, Pass Rate And Psychometric Properties Of The Test In End Of Clerkship Ophthalmology Assessment & 2022 & OA & supports the claim (0.85) \\
The Impact of the Item Types and Number of Solution Steps of Multiple-Choice Items on Item Difficulty and Discrimination and Test Reliability & 2014 & OA & supports the claim (0.40) \\
The level of Test--wiseness for the students of arts and science Faculty at Sharourah and its relationship with some variables & 2015 & OA & supports the claim (0.45) \\
The MCG Medical Teacher's Handbook & 2009 & OA & supports the claim (0.45) \\
The modified essay question: Its exit from the exit examination? & 2010 & OA & supports the claim (0.35) \\
The occurence of test-wiseness and the possibility of inducing it via instruction : a cross-sectional study & 1975 & OA & supports the claim (0.90) \\
The pearls and pitfalls of setting high-quality multiple choice questions for clinical medicine & 2023 & OA & supports the claim (0.70) \\
The Prevalence of Item Construction Flaws in Medical School Examinations and Innovative Recommendations for Improvement & 2017 & OA & supports the claim (0.80) \\
The Psychological Effect of Errors in Standardized Language Test Items on EFL Students' Responses to the Following Item & 2017 & OA & supports the claim (0.40) \\
THE QUALITY OF ISLAMIC RELIGIOUS EDUCATION TEST ITEMS IN VOCATIONAL HIGH SCHOOLS & 2026 & OA & supports the claim (0.40) \\
The Role of Faculty Development in Improving the Quality of Multiple-Choice Questions in Dental Education & 2020 & WoS, OA & supports the claim (0.55) \\
The use of guidelines for writing exam questions in medicine and the effect on the psychometric properties of the exams & 2017 & OA & supports the claim (0.55) \\
The Use of Multiple-choice Questions as an Assessment Tool in First-year University Physics Modules & 2025 & OA & supports the claim (0.80) \\
The Use of One-, Two-, and Three-Parameter and Nominal Item Response Scoring in Place of Number-Right Scoring in the Presence of Test-Wiseness & 2005 & OA & supports the claim (0.70) \\
The Validity of Two Item-Writing Rules & 1991 & OA & supports the claim (0.60) \\
Toward Evidence-Based Guidelines on Answer Keys' Placement in Multiple-Choice Tests: A Simulation Study & 2025 & OA & supports the claim (0.55) \\
Towards an authoring system for item construction & 1988 & OA & supports the claim (0.40) \\
TYPES OF ITEM WRITING FLAWS IN MULTIPLE CHOICE QUESTIONS IN MEDICAL EDUCATION AND THEIR EFFECT ON PASSING RATE OF STUDENTS & 2023 & OA & supports the claim (0.85) \\
Una vista comparativa entre los resultados obtenidos al aplicar un mismo taller de capacitación en las modalidades presencial y virtual & 2021 & OA & supports the claim (0.60) \\
USE OF AN INCLUSIVE OPTION AND THE OPTIMAL NUMBER OF OPTIONS FOR MULTIPLE-CHOICE ITEMS & 1993 & WoS, OA & supports the claim (0.80) \\
Use of Crossed Random-Effects Models to Assess Multiple-Choice Items: An Experimental Study & 2026 & OA & supports the claim (0.80) \\
Use of flawed multiple-choice items by the New England Journal of Medicine for continuing medical education & 2006 & WoS, OA & supports the claim (0.55) \\
User experience of the Written Exam Question Quality tool to inform the writing of new written-exam questions & 2024 & OA & supports the claim (0.60) \\
Using a Rule-Based System to Identify Flawed Multiple-Choice Items & 2025 & OA & supports the claim (0.50) \\
Using Test-Taking Strategies to Maximize Multiple-Choice Test Scores & 1986 & OA & supports the claim (0.75) \\
Utility of Complex Alternatives in Multiple-Choice Items: The Case of All of the Above & 2015 & OA & supports the claim (0.80) \\
Uzaktan Eğitimde Çoktan Seçmeli Soruların Güçlük ve Ayırt Edicilik Değerlerinin Soru Türlerine Göre İncelenmesi & 2016 & OA & supports the claim (0.65) \\
Validation of a Rubric to Assess Physiology Exam Multiple Choice Questions to Create Fair and Equitable Assessments & 2024 & OA & supports the claim (0.60) \\
Vermessene Fragen. Metriken als Ansatz automatisierter analytischer und konstruktiver Qualitätssicherung von Mehrfachauswahlfragen für mobile digitale Medien & 2020 & OA & supports the claim (0.70) \\
W POSZUKIWANIU TECHNOLOGII PISANIA ZADAŃ TESTÓW SPRAWDZAJĄCYCH OSIĄGNIĘCIA UCZNIÓW & 1991 & OA & supports the claim (0.55) \\
What Do Teachers Know about Measurement and How Did They Find Out & 1992 & OA & supports the claim (0.60) \\
What New And Adjunct Faculty Need To Know About Exams, Grades, And Cheating & 2011 & OA & supports the claim (0.40) \\
What Teachers Should Know about Test-Wiseness. & 1987 & OA & supports the claim (0.60) \\
When Majority Doesn't Rule: The Use of Discrimination Indices to Improve the Quality of MCQs & 2011 & OA & supports the claim (0.50) \\
When the Items Fail the Test: Prevalence of Item-writing Flaws in Multiple-choice Questions used in Introductory Forensic Science Courses & 2026 & OA & supports the claim (0.50) \\
WHICH EXAMINEES ARE MOST FAVOURED BY THE USE OF MULTIPLE CHOICE TESTS? & 1974 & OA & supports the claim (0.70) \\
Which of the Following Is True: We Can Write Better Multiple Choice Questions & 2018 & OA & supports the claim (0.70) \\
Will a Short Training Session Improve MultipleChoice ItemWriting Quality by Dental School Faculty? A Pilot Study & 2017 & OA & supports the claim (0.50) \\
Writing a multiple choice test : Do's and don'ts & 2013 & OA & supports the claim (0.40) \\
Writing a multiple-choice test question & 1992 & OA & supports the claim (0.60) \\
Writing Multiple-Choice Items & 2017 & OA & supports the claim (0.65) \\
Writing Multiple-Choice Test Items & 2023 & OA & supports the claim (0.50) \\
Writing Multiple-Choice Test Items & 1994 & OA & supports the claim (0.85) \\
Writing Multiple-Choice Test Items. ERIC/AE Digest. & 1995 & OA & supports the claim (0.50) \\
Writing multiplechoice test items & 1995 & OA & supports the claim (0.40) \\
YABANCI DİL OLARAK TÜRKÇE ÖĞRETİMİDE SINAV HAZIRLAMADA MADDE TÜRÜ YAZIMINDA DİKKAT EDİLMESİ GEREKEN DURUMLAR & 2025 & OA & supports the claim (0.65) \\
Đánh giá cht lưng b câu hi sát hch lý thuyt lái xe mô tô Hng A1 mi & 2021 & OA & supports the claim (0.70) \\
\textquotedbl{}None of the Above\textquotedbl{} as an Answer Option in Observatoin Based Multiple-Choice Questions & 2006 & OA & qualifies or contradicts the claim (0.55) \\
A Closer Look at Chinese EFL Learners' Test-wiseness Strategies in Reading Test & 2016 & OA & qualifies or contradicts the claim (0.70) \\
A METAANALYTIC REVIEW OF ITEM DISCRIMINATION AND DIFFICULTY IN MULTIPLE-CHOICE ITEMS USING NONE-OF-THE-ABOVE & 1992 & WoS, OA & qualifies or contradicts the claim (0.85) \\
A Quantitative Review of Research on Multiple-Choice Item Writing. & 1985 & OA & qualifies or contradicts the claim (0.60) \\
Adapting the Guidelines for Writing Selected-Response Items to the Scale for Acquisition of Japanese Word Meanings & 2023 & OA & qualifies or contradicts the claim (0.65) \\
AI-Enabled Quality Assurance for Multiple-Choice Assessment Items & 2026 & OA & qualifies or contradicts the claim (0.45) \\
An Empirical Evaluation of Selected Multiple-Choice Item Writing Guidelines. & 1998 & OA & qualifies or contradicts the claim (0.85) \\
Associations between item characteristics and statistical performance for paediatric medical student multiple choice assessments & 2023 & OA & qualifies or contradicts the claim (0.50) \\
Attitudes of Language Teachers toward Multiple-Choice Item Writing Guidelines: An Exploratory Factor Analysis & 2020 & OA & qualifies or contradicts the claim (0.50) \\
Comparison of Multiple-Choice Question Formats in a First Year Medical Physiology Course & 2024 & OA & qualifies or contradicts the claim (0.55) \\
Do Item Writing Best Practices Improve Multiple Choice Questions for University Students? & 2018 & OA & qualifies or contradicts the claim (0.40) \\
Do the Guideline Violations Influence Test Difficulty of High-stake Test?: An Investigation on University Entrance Examination in Turkey & 2016 & OA & qualifies or contradicts the claim (0.70) \\
Does Using None-of-the-Above (NOTA) Hurt Students' Confidence? & 2023 & OA & qualifies or contradicts the claim (0.55) \\
Editorial Changes and Item Performance: Implications for Calibration and Pretesting & 2014 & OA & qualifies or contradicts the claim (0.70) \\
Effect of faculty training on quality of multiple-choice questions & 2020 & OA & qualifies or contradicts the claim (0.45) \\
Effects of Test-Taking Strategy and Lexico-Grammatical Ability on L2 Local-Level Reading Comprehension & 2024 & OA & qualifies or contradicts the claim (0.60) \\
MULTIPLECHOICE QUESTIONS AND STUDENT CHARACTERISTICS1 & 1967 & OA & qualifies or contradicts the claim (0.60) \\
None of the Above, Less of the Right Parallel Patterns in Human and LLM Performance on Multi-Choice Questions Answering & 2025 & OA & qualifies or contradicts the claim (0.50) \\
None of the Above, Less of the Right: Parallel Patterns between Humans and LLMs on Multi-Choice Questions Answering & 2025 & OA & qualifies or contradicts the claim (0.50) \\
Randomize it: Fair Procedures when Constructing Multiple-choice Test-Keys & 2019 & OA & qualifies or contradicts the claim (0.50) \\
Scoring and Keying Multiple Choice Tests: A Case Study in Irrationality & 2004 & OA & qualifies or contradicts the claim (0.50) \\
TEST WISDOM SKILLS AND THEIR RELATIONSHIP WITH ACADEMIC ACHIEVEMENT IN HISTORY AMONG FIFTH GRADE LITERARY STUDENTS & 2025 & OA & qualifies or contradicts the claim (0.40) \\
The arrangement of response options in multiple-choice test items: verticality is not always better & 2024 & OA & qualifies or contradicts the claim (0.60) \\
The Effect of Defective Options on Multiple-choice Items and Test Characteristics & 2006 & OA & qualifies or contradicts the claim (0.70) \\
The Effective Use of Negative Stems and ``All of the Above'' in Multiple-Choice Tests in College Courses & 2020 & OA & qualifies or contradicts the claim (0.90) \\
The Effectiveness of a \textquotedbl{}None of the Above\textquotedbl{} Answer Option in Multiple Choice Grammar Tests & 2008 & OA & qualifies or contradicts the claim (0.90) \\
The impact of student ability and method for varying the position of correct answers in classroom multiple-choice tests & 2010 & OA & qualifies or contradicts the claim (0.45) \\
The Impact of the ``all-of-the-above'' Option and Student Ability on Multiple Choice Tests & 2007 & OA & qualifies or contradicts the claim (0.90) \\
The Interplay between Key Responses, Distractors, and Position Effects in Multiple-Choice Items: A Complex and Unresolved Issue & 2026 & OA & qualifies or contradicts the claim (0.40) \\
The None-of-the-Above Option: An Empirical Study & 1991 & OA & qualifies or contradicts the claim (0.90) \\
The Position of Distractors in Multiple-Choice Test Items: The Strongest Precede the Weakest & 2021 & OA & qualifies or contradicts the claim (0.35) \\
Three, four, and none of the above options in multiple-choice items & 2017 & OA & qualifies or contradicts the claim (0.70) \\
Validity of a Taxonomy of Multiple-Choice Item-Writing Rules & 1989 & OA & qualifies or contradicts the claim (0.60) \\
What Makes a Test Difficult? Exploring the Effect of Item Content on Students' Performance & 2012 & OA & qualifies or contradicts the claim (0.40) \\
\end{longtable}
\endgroup



% ---------------------------------------------------------------- claim C
% Author's notes: scholar session "rapport-parsons"; export
%   scholar sessions export rapport-parsons -f table --bearing-only \
%     --lang en --label sok-parsons -o parsons-full
% Query strings per provider: the agent's parsons-runs/ directory of
% 2026-11-13; provenance per source in ltnotes.bib.
\chapter{How are ordering items used in exams, and is one order enough?}
\label{app:parsons}
\chapterprecis{The author has not yet reviewed the results in this
  appendix in full.}

\Cref{sec:items} says that ordering items, Parsons problems, are used on
programming exams and relate to writing code, and that such an item can
have more than one valid order, which grading against one order marks
wrong; \cref{sec:checker} adds that testing every order is a way to find
them.  The questions are:
\begin{description}
  \item[(a)] Are Parsons problems used in programming exams, and do their
    scores relate to performance on writing code?
  \item[(b)] Does the literature recognize that a Parsons problem can have
    several valid orders, so that grading against one fixed order marks a
    correct answer wrong?
\end{description}
The report needs (a) at the strength \enquote{used, and related}, not
\enquote{measuring the same skill}, and (b) as a recognized problem.

\section{Method}
\label{sec:parsons-method}

The searches ran on 13 and 14 November 2026 with \texttt{scholar}, in the
same databases and in the same way as in \cref{sec:item-flaws-method}; one
OpenAlex counter-query failed on the shared budget and ran again after it
was reset.  All queries share a concept block P with the field's names for
the item type.  Three support queries add exams (S1), writing or tracing
code (S2), and grading or several solutions (S3); two counter-queries add
words for criticism and weak correlation to the terms of S1 and S2 (M1)
and of S3 (M2).  \Cref{tab:sok-queries-parsons} gives them with their
hits.

\begin{table}[htbp]
  \caption{Queries and hits per database for the claims of
    \cref{app:parsons}, 13--14 November 2026, tool \texttt{scholar}, no
    limit on hits.  P stands for (\enquote{Parsons problem} OR
    \enquote{Parsons problems} OR \enquote{Parson's problem} OR
    \enquote{Parson's problems} OR \enquote{Parsons puzzle} OR
    \enquote{Parsons puzzles} OR \enquote{Parson's puzzle} OR
    \enquote{Parsons programming puzzle} OR \enquote{Parsons programming
    puzzles} OR \enquote{mixed-up code} OR \enquote{code scramble} OR
    \enquote{code ordering} OR \enquote{drag-and-drop programming});
    quoted words were searched as phrases.  S = support query,
    M = counter-query.  Databases as in \cref{tab:sok-queries-item-flaws}:
    Web of Science (WoS), Scopus, IEEE Xplore (IEEE), OpenAlex (OA, full
    text).  No search failed in the end.}
  \label{tab:sok-queries-parsons}
  \centering\footnotesize
  \begin{tabularx}{\linewidth}{@{}l>{\raggedright\arraybackslash}X
    rrrr@{}}
    \toprule
       & Keywords & WoS & Scopus & IEEE & OA \\
    \midrule
    S1 & P AND (exam OR exams OR examination OR assessment OR test)
       & 46 & 75 & 9 & 3\,358 \\
    S2 & P AND (\enquote{code writing} OR \enquote{writing code} OR
         \enquote{code-writing} OR tracing OR correlation OR correlate)
       & 32 & 51 & 3 & 1\,898 \\
    S3 & P AND (grading OR scoring OR \enquote{multiple solutions} OR
         \enquote{alternative solutions} OR \enquote{multiple correct} OR
         \enquote{unique solution} OR \enquote{execution-based} OR
         \enquote{line-based})
       & 26 & 10 & 6 & 1\,772 \\
    \midrule
    M1 & P AND (the terms of S1 OR the terms of S2) AND (limitation OR
         limitations OR critique OR criticism OR validity OR \enquote{weak
         correlation} OR \enquote{poor correlation} OR \enquote{no
         correlation} OR \enquote{not correlate})
       & 1 & 4 & 0 & 2\,770 \\
    M2 & P AND (the terms of S3) AND (limitation OR limitations OR
         critique OR criticism OR validity OR reliability OR fairness OR
         \enquote{not a problem})
       & 0 & 2 & 0 & 1\,675 \\
    \bottomrule
  \end{tabularx}
\end{table}

The searches gave 3\,635 unique records, almost all from OpenAlex's
full-text search, where \enquote{code ordering} and \enquote{test} occur in
many unrelated texts.  A language model classified every record as in
\cref{sec:item-flaws-method}, with the discard reasons another field, no
Parsons problem, Parsons problems as practice only, and no evidence on
either question.  We read every kept record, every record the model was
unsure of (75), a sample of 25 discards from other fields, and every hit
of the counter-queries in Web of Science, Scopus and IEEE Xplore.  Of the
11 cited sources, six were read in full text.  Five, all in the ACM
Digital Library, were read only in their abstracts, because the library
blocked automated access; among them is the first study of Parsons
problems as exam questions.

\section{Results}
\label{sec:parsons-results}

\paragraph{Used in exams.}
A working-group review of the field notes that
\textcquote[207]{ericson2022}{paper-based delivery has continued to be
popular for use in exams.  In 2008, the first paper that explored the use
of Parsons problems on exams was published.}  Since then they have been
used on computer-based exams in introductory programming in Python: one
course \textcquote[674]{wu2024}{included a variety of micro Parsons
problems on four exams}; another study gives item statistics for
\textcquote[417]{smith2023}{32 pairs of Parsons problems from CS1 Python
exams and quizzes}; a third assessed code sequencing with Parsons problems
on four exams\autocite[4]{newar2025}.

\paragraph{Related to writing code, but not the same.}
The first exam study reports, in its abstract,
\textcquote[113]{denny2008}{notable correlation between Parsons scores and
code writing scores} --- and \enquote{low correlation} between Parsons
problems and tracing code.  Later work qualifies the relation.  One-line
Parsons items were \textcquote[679]{wu2024}{comparable in difficulty but
slightly lower in discrimination compared to writing code from scratch};
a study of four exams \textcquote{newar2025}{did not find a clear
prerequisite relationship between sequencing code and writing or
explaining code}, where an earlier one had placed Parsons puzzles at an
intermediate level below writing code\autocite{lopez2008}; and a review
finds the effectiveness of Parsons problems \textcquote[195]{du2020}{both
as a question type and as a learning tool in CS education} uncertain
\enquote{due to a lack of replicated research}.

\paragraph{Several valid orders.} The problem is named early:
\textcquote[IIP-127]{ihantola2011}{Ambiguous solutions are perhaps the
biggest problem of automatic assessment of complex Parson's puzzles.
\textelp{} This is a problem if the assessment of Parson's puzzles is
based on the correct order of the lines.}  The same authors propose what
the rewrite did: puzzles \textcquote[IIP-127]{ihantola2011}{should be
constructed in an environment that tests all the possible combinations,
alerts the author about alternative solutions, and leaves it for the
author to decide which of the solutions should be accepted}.  The review
states the trade-off: \textcquote[218]{ericson2022}{execution-based
feedback can allow multiple possible solutions, while line-based systems
typically have only one correct solution}.  For ordering items in proofs,
the instructor can give the dependency graph of the lines, so that
\textcquote[428]{poulsen2022}{any correct arrangement of the lines can
receive full credit}; its authors report from practice that when an
instructor misses a valid rearrangement, the result is
\textcquote[431]{poulsen2022}{a correct student response being incorrectly
marked as faulty by the autograder}.  A recent preprint adds why
order-based grading persists: instructors prefer it because it
\textcquote{oakeson2025}{can give students more information about how to
fix their solutions, does not require writing test cases}.

\section{Discussion and limitations}
\label{sec:parsons-discussion}

The databases' limits are those of \cref{sec:item-flaws-discussion}, and
the conference proceedings of computing education are indexed mainly
through OpenAlex and Scopus.  The most important limitation is that the
correlation with writing code rests on one study read only in its
abstract: its size, sample and grading are unverified.  The most direct
report of correct answers marked wrong concerns proofs, not programs; for
programs the evidence is a design argument with an example.  Ihantola and
Karavirta also qualify (b) themselves: some alternative orders, such as
ones that put distractors where they never run, are bad programming and
should be marked wrong.

\section{Conclusion}
\label{sec:parsons-conclusion}

Yes to both, with qualifications.  Parsons problems are used on
programming exams, including computer-based exams in Python, and their
scores correlate notably with writing code in the first study of them,
though later work shows that they are not interchangeable with writing
code.  That an ordering item can have several valid orders is recognized
as a central problem of grading by order, and testing every order, the
remedy the rewrite used, is the one its first describers proposed.  The
report therefore says \enquote{used and related}, not \enquote{the same
skill}, and treats a second valid order as a defect to find, deciding for
each found order whether it should count.

\section{Further work}
\label{sec:parsons-further}

The five sources read only in abstract should be read in full, the first
exam study first.\footnote{Denny, Luxton-Reilly and Simon 2008,
\texttt{doi:10.1145/1404520.1404532}.}  No study found measures how often
ordering items on real exams have alternative valid orders, or how many
students lose points through them; regrading an exam's ordering items with
test-based grading and counting the changed outcomes would settle it.  For
(a), a comparison of the same tasks as Parsons problems and as code
writing, with whole programs, in an introductory Python exam would show
whether one can stand in for the other.

\section{Hits bearing on the claim}
\label{sec:parsons-hits}

\Cref{tab:sok-parsons} lists every record the screening kept as cited,
supporting or qualifying.
% --- scholar LaTeX fragment --------------------------------------------
% Generated by: scholar sessions export -f table
% Session: rapport-parsons
% \input this file from the body of a document whose
% preamble loads:
%   \usepackage{longtable}
%   \usepackage{booktabs}
%   \usepackage{array}
% ----------------------------------------------------------------------

\begingroup\footnotesize
\begin{longtable}{@{}%
  >{\raggedright\arraybackslash}p{0.44\textwidth}c%
  >{\raggedright\arraybackslash}p{0.13\textwidth}%
  >{\raggedright\arraybackslash}p{0.26\textwidth}@{}}
\caption{Records bearing on the claim, Parsons problems as exam items (11 cited, 18 supporting and 4 qualifying or contradicting it, out of 3635 unique hits; the 618 adjacent and 2984 false hits are not listed). The last column gives the reason for each inclusion; for machine-classified rows the model's confidence is given. Databases: OpenAlex (OA; not open access), IEEE Xplore (IEEE), Web of Science (WoS), Scopus.}\label{tab:sok-parsons}\\
\toprule
Title & Year & Provider & Reason \\
\midrule
\endfirsthead
\multicolumn{4}{@{}l}{\emph{\tablename~\thetable{} (cont.)}}\\
\toprule
Title & Year & Provider & Reason \\
\midrule
\endhead
\midrule \multicolumn{4}{r@{}}{\emph{continued on next page}}\\
\endfoot
\bottomrule
\endlastfoot
A Review of Research on Parsons Problems & 2020 & OA & \emph{cited}: review \\
Choose Your Own Solution: Supporting Optional Blocks in Block Ordering Problems & 2025 & OA & \emph{cited}: any-valid-order grading \\
Evaluating a new exam question: Parsons problems & 2008 & OA & \emph{cited}: Parsons vs code writing \\
Evaluating Micro Parsons Problems as Exam Questions & 2024 & OA & \emph{cited}: Parsons items on exams \\
Evaluating Proof Blocks Problems as Exam Questions & 2021 & OA & \emph{cited}: Parsons items on exams \\
Investigating the Role and Impact of Distractors on Parsons Problems in CS1 Assessments & 2023 & OA & \emph{cited}: Parsons items on exams \\
Mining hierarchies with conviction: constructing the CS1 skill hierarchy with pairwise comparisons over skill distributions & 2025 & OA & \emph{cited}: skill hierarchy \\
Parsons Problems and Beyond: Systematic Literature Review and Empirical Study Designs & 2022 & OA & \emph{cited}: review \\
Proof Blocks: Autogradable Scaffolding Activities for Learning to Write Proofs & 2022 & OA & \emph{cited}: any-valid-order grading \\
Relationships between reading, tracing and writing skills in introductory programming & 2008 & OA & \emph{cited}: skill hierarchy \\
Two-Dimensional Parson's Puzzles: The Concept, Tools, and First Observations & 2011 & OA & \emph{cited}: multiple valid orders \\
Analysis of the Answering Processes in Split-Paper Testing to Promote Instruction & 2021 & OA & supports the claim (0.60) \\
Benchmarking Partial Credit Grading Algorithms for Proof Blocks Problems & 2022 & OA & supports the claim (0.90) \\
Beyond MCQ: Designing Engaging, Autogradable Assessments for Supporting Teaching \& Learning in K-12 Computer Science & 2022 & OA & supports the claim (0.50) \\
Comprehension analysis considering programming thinking ability using code puzzle & 2020 & OA & supports the claim (0.50) \\
Detecting Struggling Students from Interactive Ebook Data & 2022 & OA & supports the claim (0.65) \\
Efficient Feedback and Partial Credit Grading for Proof Blocks Problems & 2023 & OA & supports the claim (0.90) \\
Epplets: A Tool for Solving Parsons Puzzles & 2018 & OA & supports the claim (0.55) \\
Extension of a Study on Parsons' Problems Reveals Implications for SEE\&T & 2024 & OA & supports the claim (0.85) \\
Generating Multi-Part Autogradable Faded Parsons Problems From Code-Writing Exercises & 2024 & OA & supports the claim (0.55) \\
How Do Students Solve Parsons Programming Problems? -- Execution-Based vs. Line-Based Feedback & 2013 & OA & supports the claim (0.50) \\
Integrating Parsons problems in a CS1 programming course autograder & 2024 & OA & supports the claim (0.40) \\
Open Source Widget for Parson's Puzzles & 2010 & WoS, Scopus & supports the claim (0.60) \\
Parson's Puzzles & 2013 & OA & supports the claim (0.60) \\
Process-Oriented Understanding Estimation Using Code Puzzles & 2022 & OA & supports the claim (0.70) \\
Redefining Evaluation Using Adaptive AI-Generated Parsons Problems & 2026 & OA & supports the claim (0.60) \\
Split-Paper Testing: A Novel Approach to Evaluate Programming Performance & 2020 & OA & supports the claim (0.90) \\
The Impact of Solving Adaptive Parsons Problems with Common and Uncommon Solutions & 2022 & OA & supports the claim (0.50) \\
Using a Computer to Score Parsons Problems Answered on Paper & 2021 & OA & supports the claim (0.90) \\
Comparing the Impacts of Visually Grouped and Jumbled Distractors on Parsons Problems in CS1 Assessments & 2023 & OA & qualifies or contradicts the claim (0.65) \\
Discovering, autogenerating, and evaluating distractors for Parsons problems in CS1 & 2025 & OA & qualifies or contradicts the claim (0.60) \\
Discovering, Autogenerating, and Evaluating Distractors for Python Parsons Problems in CS1 & 2023 & OA & qualifies or contradicts the claim (0.80) \\
Useful Distractions? Investigating the Utility of Distractors in Parsons Problems & 2023 & OA & qualifies or contradicts the claim (0.60) \\
\end{longtable}
\endgroup



% ---------------------------------------------------------------- claim D
% Author's notes: scholar session "rapport-feedback" (an earlier combined
% run set aside as "rapport-feedback-combined-run0"); export
%   scholar sessions export rapport-feedback -f table --bearing-only \
%     --lang en --label sok-feedback -o feedback-full
% Provenance per source in ltnotes.bib.
\chapter{Does feedback that explains help more than right or wrong?}
\label{app:feedback}
\chapterprecis{The author has not yet reviewed the results in this
  appendix in full.}

\Cref{sec:workflow} says that every item got feedback that explains ---
what to think about, how to check it, where to study --- and that such
feedback helps learning more than feedback that only says right or wrong,
though not always.  The questions are:
\begin{description}
  \item[(a)] Does elaborated feedback on test items, which explains why an
    answer is right or wrong or what to study, support learning more than
    verification feedback, which says right or wrong or shows the correct
    answer, in particular on computer-based items?
  \item[(b)] Does feedback sometimes fail to help, or even lower
    performance?
\end{description}
The report needs (a) on average, with its limits, and (b) as a caution.

\section{Method}
\label{sec:feedback-method}

The searches ran on 13 and 14 November 2026 with \texttt{scholar}, in the
same databases and in the same way as in \cref{sec:item-flaws-method};
the five OpenAlex searches ran after its daily budget was reset.  The
literature on feedback is large, so the queries combine narrow concept
blocks rather than single words: elaborated feedback (E), verification
feedback under the field's names for it (V), computer-based or quiz items
(C), reviews and meta-analyses (R), and feedback that lowers performance.
Two counter-queries add to E and V words for no difference (M1) and for
learners ignoring feedback (M2).  \Cref{tab:sok-queries-feedback} gives
them with their hits.

\begin{table}[htbp]
  \caption{Queries and hits per database for the claims of
    \cref{app:feedback}, 13--14 November 2026, tool \texttt{scholar}, no
    limit on hits.  E stands for (\enquote{elaborated feedback} OR
    \enquote{elaborative feedback} OR \enquote{explanatory feedback}),
    V for (\enquote{verification feedback} OR \enquote{knowledge of
    results} OR \enquote{knowledge of correct response} OR
    \enquote{corrective feedback}); quoted words were searched as
    phrases.  S = support query, M = counter-query.  Databases as in
    \cref{tab:sok-queries-item-flaws}: Web of Science (WoS), Scopus, IEEE
    Xplore (IEEE), OpenAlex (OA, full text).  No search failed in the
    end.}
  \label{tab:sok-queries-feedback}
  \centering\footnotesize
  \begin{tabularx}{\linewidth}{@{}l>{\raggedright\arraybackslash}X
    rrrr@{}}
    \toprule
       & Keywords & WoS & Scopus & IEEE & OA \\
    \midrule
    S1 & E AND V AND (\enquote{computer-based} OR \enquote{multiple-choice}
         OR quiz OR quizzes OR \enquote{test items})
       & 19 & 15 & 2 & 709 \\
    S2 & E AND V AND (\enquote{meta-analysis} OR \enquote{meta-analytic} OR
         review)
       & 10 & 7 & 1 & 1\,013 \\
    S3 & (\enquote{feedback intervention} OR \enquote{feedback
         interventions} OR \enquote{effects of feedback} OR
         \enquote{formative feedback}) AND (\enquote{meta-analysis} OR
         \enquote{meta-analytic}) AND (\enquote{negative effect} OR
         \enquote{negative effects} OR \enquote{decrease performance} OR
         \enquote{lower performance} OR debilitating) AND (E OR V)
       & 0 & 0 & 0 & 692 \\
    \midrule
    M1 & E AND V AND (\enquote{no difference} OR \enquote{no significant
         difference} OR \enquote{not more effective} OR ineffective OR
         \enquote{small effect} OR inconsistent OR \enquote{no effect})
       & 1 & 5 & 0 & 1\,029 \\
    M2 & E AND V AND (ignore OR ignored OR \enquote{attention to feedback}
         OR \enquote{feedback use} OR \enquote{use of feedback} OR
         \enquote{time spent on feedback})
       & 1 & 2 & 0 & 574 \\
    \bottomrule
  \end{tabularx}
\end{table}

The searches gave 1\,648 unique records.  A language model classified
every one as in \cref{sec:item-flaws-method}, with the discard reasons
correction of second-language writing, motor learning, no comparison of
feedback types, description of a tool only, and not about feedback.  We
read the qualifying records closest to the claim, about a dozen of the 77,
and decided the cited ones by hand.  One review of formative feedback, the
standard reference for the distinction between verification and
elaboration, was not returned by the searches and was retrieved as a known
document by its identifier, since the author had named it as a source to
check. Six of the nine cited sources were read in full text; one
meta-analysis in its submitted version, as a chapter of its first author's
thesis, so its pages are the thesis's. Three were read only in their
abstracts.  Two classic reviews of feedback were not returned by these
searches and could not be read in full, so they are not cited; one of them
is reported here through a later review.

\section{Results}
\label{sec:feedback-results}

\paragraph{Elaborated against verification feedback.}
A meta-analysis of feedback on items in computer-based environments found
that \textcquote{vanderkleij2015}{elaborated feedback (EF; e.g., providing
an explanation) produced larger effect sizes (0.49) than feedback regarding
the correctness of the answer (KR; 0.05) or providing the correct answer
(KCR; 0.32)}.  A later network meta-analysis agrees that
\textcquote[1743]{mertens2022}{EF is most likely to be the most effective
for lower-order (i.e., recall/recognition) and higher-order (i.e.,
transfer) learning outcomes compared with the other feedback variants},
and a review of formative feedback that it is more effective
\textcquote[157]{shute2008}{when it provides details of how to improve the
answer rather than just indicating whether the student's work is correct
or not}.

The advantage has limits.  With multiple-choice items and lower-order
outcomes, \textcquote[1764]{mertens2022}{EF, AUC, and KCR feedback were
all equally effective}, AUC being answer-until-correct; and that analysis
covered only \textcquote[1765]{mertens2022}{the effects of immediate
computer-based feedback}.  In a randomized experiment in higher education
\textcquote[23]{vanderkleij2012}{\textins*{n}o significant effect was
found for the feedback condition on student achievement on the post-test};
in a quiz app \textcquote[1]{rueth2021}{cognitive outcomes (quiz scores)
and metacognitive outcomes (response certainty) increased similarly}
whatever the feedback type; after computer-based exams,
\textcquote{huebner2024}{\textins*{c}ompared with EF, KCR provided modest
improvements on follow-up exams, specifically on repeated questions}; and
in a primary school \textcquote{caceres2021}{the Knowledge of Results
group scored higher on the post-test}.  The last two were read only in
their abstracts and are small, quasi-experimental studies.

\paragraph{Whether students use it.}
Students \textcquote[38]{vanderkleij2012}{spent more time reading feedback
when the feedback was delivered immediately than when the feedback was
delivered with a delay}, and \textcquote[65]{manez2020}{\textins*{s}tudents who
received KCR feedback seldom access EF}, although that study found
\textcquote[57]{manez2020}{no differences in the final learning}.
Elaborate feedback is no guarantee either: \textcquote[159]{shute2008}{if
feedback is too long or too complicated, many learners will simply not pay
attention to it}.  On timing the evidence is mixed: delayed feedback
weakened the effects in the first meta-analysis\autocite{vanderkleij2015},
while in a small study of multiple-choice tests \textcquote{ryan2024}{the
feedback timing effect was non-significant}.

\paragraph{Feedback that does not help.}
Feedback can fail outright: \textcquote[156]{shute2008}{about one third of
the total studies reviewed in two landmark meta-analyses \textelp{}
demonstrate negative effects of feedback on learning}, one of them a
meta-analysis of feedback interventions in general that we report only
through this review.  In the meta-analysis of computer-based item feedback,
\textcquote[77]{vanderkleij2015}{the number of effect sizes that were
negative was 12, and for each feedback type, there were both negative and
positive effect sizes}.

\section{Discussion and limitations}
\label{sec:feedback-discussion}

The databases' limits are those of \cref{sec:item-flaws-discussion}; here
the title-and-abstract databases found little, and almost all records came
from OpenAlex's full-text search.  The model's classification recorded no
confidence, and we read only the qualifying records closest to the claim.
The page numbers of one meta-analysis are those of its submitted version.
Three studies were read only in abstract.  The two classic reviews the
field cites most were neither returned by the searches nor obtainable in
full.  Most evidence comes from higher education, little of it from
programming, and almost none concerns feedback that students read after a
summative exam.

\section{Conclusion}
\label{sec:feedback-conclusion}

Yes to (a) on average, with limits, and yes to (b).  Meta-analyses of
feedback on computer-based items find that feedback that explains helps
learning clearly more than feedback that only says right or wrong, and
somewhat more than showing the correct answer; but the margin over the
correct answer is small, absent for multiple-choice items in one
analysis, and single studies find no difference or the reverse.  Feedback
does not always help: some effects are negative, and students often skip
it, especially when it comes late.  For the report this means: the
explaining feedback is defensible against right-or-wrong feedback, not
clearly better than showing the answer, and best supported where it comes
right after an attempt, as in a practice exam.

\section{Further work}
\label{sec:feedback-further}

This search left the published versions of two sources read in other
versions, the two classic reviews, and one study of elaboration and timing
in computer-based practice to be read; the last could not be read because
the copy behind its identifier was another
paper.\footnote{Attali and van der Kleij 2017, \emph{Computers \&
Education}, \texttt{doi:10.1016/j.compedu.2017.03.012}.}  We read only
about a dozen of the 77 qualifying records.

The field lacks studies of written feedback on multiple-choice items that
students read after a summative computer exam, and of programming in
particular.  A randomized study that compares, in practice and after an
exam, showing the correct answer with an explanation that does not name
it, and logs how long students read, would show whether texts like these
are worth writing; if they were not, the effort would be better spent on a
practice exam.

\section{Hits bearing on the claim}
\label{sec:feedback-hits}

\Cref{tab:sok-feedback} lists every record the screening kept as cited,
supporting or qualifying.
% --- scholar LaTeX fragment --------------------------------------------
% Generated by: scholar sessions export -f table
% Session: rapport-feedback
% \input this file from the body of a document whose
% preamble loads:
%   \usepackage{longtable}
%   \usepackage{booktabs}
%   \usepackage{array}
% ----------------------------------------------------------------------

\begingroup\footnotesize
\begin{longtable}{@{}%
  >{\raggedright\arraybackslash}p{0.44\textwidth}c%
  >{\raggedright\arraybackslash}p{0.13\textwidth}%
  >{\raggedright\arraybackslash}p{0.26\textwidth}@{}}
\caption{Records bearing on the claim (8 cited, 216 supporting and 77 qualifying or contradicting it, out of 1648 unique hits; the 639 adjacent and 708 false hits are not listed). The last column gives the reason for each inclusion; for machine-classified rows the model's confidence is given. Databases: OpenAlex (OA; not open access), IEEE Xplore (IEEE), Web of Science (WoS), Scopus.}\label{tab:sok-feedback}\\
\toprule
Title & Year & Provider & Reason \\
\midrule
\endfirsthead
\multicolumn{4}{@{}l}{\emph{\tablename~\thetable{} (cont.)}}\\
\toprule
Title & Year & Provider & Reason \\
\midrule
\endhead
\midrule \multicolumn{4}{r@{}}{\emph{continued on next page}}\\
\endfoot
\bottomrule
\endlastfoot
Effects of Computer-Based Feedback on Lower and Higher-Order Learning Outcomes: A Network Meta-Analysis & 2022 & Scopus & \emph{cited}: advantage of elaborated feedback \\
Effects of feedback in a computer-based assessment for learning & 2011 & OA & \emph{cited}: whether learners use feedback \\
Effects of Feedback in a Computer-Based Learning Environment on Students' Learning Outcomes: A Meta-Analysis & 2015 & WoS & \emph{cited}: advantage of elaborated feedback \\
Exam Review Versus Categorical Feedback: Which Improves Scores by Doctor of Physical Therapy Students? & 2023 & OA & \emph{cited}: limits of elaborated feedback \\
Is more detailed feedback better for problem-solving? & 2019 & WoS & \emph{cited}: limits of elaborated feedback \\
The Effects of Different Feedback Types on Learning With Mobile Quiz Apps & 2021 & OA & \emph{cited}: limits of elaborated feedback \\
Timing's not everything: Immediate and delayed feedback are equally beneficial for performance in formative multiplechoice testing & 2023 & OA & \emph{cited}: feedback timing \\
¿Influye la Retroalimentación Correctiva en el Uso de la Retroalimentación Elaborada en un Entorno Digital? & 2019 & OA & \emph{cited}: whether learners use feedback \\
 & 2017 & OA & supports the claim (0.75) \\
  & 2024 & OA & supports the claim (0.70) \\
                & 2023 & OA & supports the claim (0.75) \\
                  & 2020 & OA & supports the claim (0.70) \\
                       & 2021 & OA & supports the claim (0.65) \\
                           & 2020 & OA & supports the claim (0.60) \\
     (  \&  )   (  \&  )                      & 2018 & OA & supports the claim (0.40) \\
     ( / )        (/ )           & 2021 & OA & supports the claim (0.60) \\
     ( / )     ( / )                & 2021 & OA & supports the claim (0.70) \\
     (/ )      (/ )           & 2023 & OA & supports the claim (0.75) \\
    ( )                    & 2020 & OA & supports the claim (0.60) \\
    (/ )        & 2020 & OA & supports the claim (0.60) \\
   ( )                         & 2019 & OA & supports the claim (0.55) \\
   ( --  )                       & 2022 & OA & supports the claim (0.50) \\
   ( -- )                     & 2019 & OA & supports the claim (0.55) \\
   (- )             & 2022 & OA & supports the claim (0.85) \\
   (/ )                       Feedback level ( corrective / explanatory ) in the learning environment based on electronic maps and its impact on the chemical and scientific & 2016 & OA & supports the claim (0.60) \\
   (/ )        (/ )            & 2020 & OA & supports the claim (0.50) \\
   4 & 2017 & OA & supports the claim (0.85) \\
 4 & 2020 & OA & supports the claim (0.80) \\
A CASE STUDY OF FEEDBACK-RICH ONLINE FORMATIVE ASSESSMENT IN BEGINNER KOREAN GRAMMAR LEARNING & 2026 & OA & supports the claim (0.40) \\
A developmental perspective on feedback: How corrective feedback influences children's literacy, mathematics, and problem solving & 2022 & OA & supports the claim (0.80) \\
A Digital Math Game and Multiple-Try Use with Primary Students: A Sex Analysis on Motivation and Learning & 2024 & OA & supports the claim (0.50) \\
A Gender and Feedback Effects in Digital Game-Based Learning for Primary Mathematics Education & 2025 & OA & supports the claim (0.60) \\
A memory advantage for past-oriented over future-oriented performance feedback. & 2018 & OA & supports the claim (0.40) \\
A Meta-Analysis of Intelligent Homework Systems (IHS): Effects on Academic Performance and the Role of Functional Features & 2026 & OA & supports the claim (0.55) \\
A review of the characteristics of intermediate assessment and their relationship with student grades & 2017 & OA & supports the claim (0.60) \\
A review of the impact of formative assessment on student achievement in mathematics & 2017 & OA & supports the claim (0.50) \\
A study of computer-assisted instructional strategies and learner characteristics & 2011 & OA & supports the claim (0.80) \\
A study of the effect of question feedback types on learning engagement in panoramic videos & 2025 & OA & supports the claim (0.65) \\
A systematic evaluation of the immediate feedback assessment technique & 2011 & OA & supports the claim (0.60) \\
A SYSTEMATIC REVIEW OF FEEDBACK MECHANISMS IN DISTANCE EDUCATION: EXPLORING SOURCES, TYPES, AIMS, FORMS, AND OUTCOMES & 2026 & OA & supports the claim (0.45) \\
Acquiring complex concepts through classification versus observation & 2024 & OA & supports the claim (0.85) \\
Adaptive feedback and student behaviour in computerassisted instruction & 2010 & OA & supports the claim (0.70) \\
Adaptives Feedback zur Unterstützung des selbstregulierten Lernens durch Experimentieren & 2011 & OA & supports the claim (0.40) \\
Affective-motivational effects of performance feedback in computer-based assessment: Does error message complexity matter? & 2022 & OA & supports the claim (0.50) \\
AI-Facilitated Formative Feedback and Students' Classroom Engagement: A Review of Pedagogical Integration, Learning Outcomes, and Teacher-Student Interaction Dynamics & 2026 & OA & supports the claim (0.55) \\
An empirical evaluation of the formative feedback supported by dashboard in the context of compilation error & 2023 & OA & supports the claim (0.55) \\
An examination of the effects of feedback type and regulation style on test achievement and review strategy using computer web-based anatomy labs with first year kinesiology students & 2000 & OA & supports the claim (0.60) \\
An investigation of using elaborated and metacognitive feedback strategies in interactive instructional videos & 2023 & OA & supports the claim (0.45) \\
Assessing the Impact and Underlying Pathways of Sequenced AI feedback on Student Learning & 2026 & OA & supports the claim (0.60) \\
Audience response systems for formative assessment purposes: The impact of individualised feedback on EFL student teachers' receptive skills & 2026 & OA & supports the claim (0.50) \\
Automated feedback and writing: a multi-level meta-analysis of effects on students' performance & 2023 & OA & supports the claim (0.50) \\
Automated Feedback Can Improve Hypothesis Quality & 2019 & OA & supports the claim (0.60) \\
Automated intelligent feedback in university statistics education & 2020 & OA & supports the claim (0.50) \\
Balancing error-correction hints and solution guidance: a quasi-experimental study of ChatGPT-integrated feedback strategies in Jupyter-based programming education & 2026 & OA & supports the claim (0.80) \\
Bedeutung von Elaboration und Kontrastierung von Falschantworten für die Entwicklung differentialdiagnostischer und -therapeutischer Entscheidungskompetenz im Medizinstudium & 2023 & OA & supports the claim (0.70) \\
BİLGİSAYAR DESTEKLİ EĞİTİM YAZILIMLARINDA KULLANILAN GERİ BİLDİRİM TÜRLERİNİN AKADEMİK BAŞARIYA ETKİSİ & 2011 & OA & supports the claim (0.80) \\
Chapter 11 Students' interaction with and appreciation of automated informative tutoring feedback & 2025 & OA & supports the claim (0.60) \\
Characteristics of Feedback that Influence Student Confidence and Performance during Mathematical Modeling & 2015 & OA & supports the claim (0.50) \\
Children's Identification and Filling of Knowledge Gaps in Digital Learning Environments & 2025 & OA & supports the claim (0.40) \\
Cognitive perspectives on maintaining physicians' medical expertise: IV. Best practices and open questions in using testing to enhance learning and retention & 2023 & OA & supports the claim (0.60) \\
Computer-based feedback in formative assessment & 2013 & OA & supports the claim (0.80) \\
Computer-based knowledge of results feedback in different delivery modes: Effects on performance, motivation, and achievement emotions & 2021 & OA & supports the claim (0.50) \\
Computer-based performance feedback: Effects of error message complexity on cognitive, metacognitive, and motivational outcomes & 2023 & OA & supports the claim (0.60) \\
Conditions for Effective Learning Without Upfront Instruction: How Practice with Feedback Supports Memory, Generalization, Motivation, and Metacognition & 2025 & OA & supports the claim (0.80) \\
CONSTRUCTIVE FEEDBACK AND STUDENTS' ACADEMIC ACHIEVEMENT: A THEORETICAL FRAMEWORK & 2020 & OA & supports the claim (0.40) \\
Constructive Feedback Intervention for Students' Academic Achievement in Chemistry: A Case of Public Secondary Schools of Pakistan & 2023 & OA & supports the claim (0.45) \\
Deliverable 4.4 - Review effective feedback and formative assessment in e-learning & 2014 & OA & supports the claim (0.70) \\
Designing automated guidance for concept diagrams in inquiry instruction & 2016 & OA & supports the claim (0.70) \\
Differential Effects of Social and Digital Feedback Interventions in the Figural Analogies Learning Test (FALT) & 2026 & OA & supports the claim (0.50) \\
Differential susceptibility in education : interaction between genes, regulatory skills, and computer games & 2011 & OA & supports the claim (0.40) \\
Digital Feedback Styles (Corrective vs Explanatory): Its Effect on Developing Secondary Students' Cognitive Achievement and Self-regulated Learning through a Digital learning Platform. & 2022 & OA & supports the claim (0.60) \\
Does Computer-based Elaborated Feedback Influence the Students' Question-Answering Process? & 2019 & OA & supports the claim (0.80) \\
Does feedback targeting text comprehension trigger the use of reading strategies or changes in readers' attitudes? A metaanalysis & 2022 & OA & supports the claim (0.75) \\
Does Student Learning Rate Depend on Feedback Type and Prior Knowledge? & 2025 & OA & supports the claim (0.80) \\
Eects of digital formative assessment tools on teaching quality and student achievement & 2020 & OA & supports the claim (0.40) \\
EFFECT OF DESCRIPTIVE FEEDBACK AND CORRECTIVE FEEDBACK ON ACADEMIC ACHIEVEMENT OF VII GRADERS IN MATHEMATICS & 2013 & OA & supports the claim (0.75) \\
Effect of Feedback and Strategy Training on Undergraduate Students' Writing Ability & 2023 & OA & supports the claim (0.55) \\
Effect of Feedback Specificity On Students' Academic Self-Efficacy at Intermediate Level & 2022 & OA & supports the claim (0.40) \\
Effect of Feedback Strategy and Motivation of Achievement to Improving Learning Results Concept in Learning Civic Education in Vocational High School & 2017 & OA & supports the claim (0.55) \\
Effect of feedback type on enhancing subsequent memory: Interaction with initial correctness and confidence level & 2021 & OA & supports the claim (0.85) \\
Effect of Immediate Elaborated Feedback on Rater Accuracy & 2020 & OA & supports the claim (0.60) \\
Effectiveness of Scaffolding in Simulated IT Training and Education & 2015 & OA & supports the claim (0.60) \\
Effects of a computer-assisted formative assessment intervention based on multiple-tier diagnostic items and different feedback types & 2016 & OA & supports the claim (0.80) \\
Effects of different AI-driven Chatbot feedback on learning outcomes and brain activity & 2025 & OA & supports the claim (0.70) \\
Effects of Different Feedback Types and Answer Changing on Achievement of Tenth Grade Students ( 4) & 2020 & OA & supports the claim (0.85) \\
Effects of Digital Score Reporting and Feedback on Students' Learning in Higher Education & 2019 & OA & supports the claim (0.60) \\
Effects of error correction strategies on long division performance of students with learning problems & 1993 & OA & supports the claim (0.45) \\
Effects of feedback elaboration and feedback timing during computer-based practice in mathematics problem solving & 2017 & OA & supports the claim (0.70) \\
Effects of Feedback on Learning Strategies in Learning Journals: Learner-Expertise Matters & 2011 & OA & supports the claim (0.70) \\
Effects of guiding feedback on students' performance, calibration, and self-efficacy: Insights from a field study with engineering students & 2026 & OA & supports the claim (0.60) \\
Effects of informative feedback on the regulation of achievement goals & 2013 & OA & supports the claim (0.40) \\
Effects of Practice Testing with Corrective Feedback on Long-Term BPSD Knowledge in Dementia Caregivers: Randomized Controlled Trial & 2026 & OA & supports the claim (0.75) \\
Effects of Using Immersive Virtual Reality for Science Education on Learning Outcomes: A Randomized Controlled Pilot Study & 2023 & OA & supports the claim (0.60) \\
Efficacité des feed-back en apprentissage numérique : effets du contenu et de la forme & 2025 & OA & supports the claim (0.85) \\
Elaborated feedback and learning: Examining cognitive and motivational influences & 2019 & OA & supports the claim (0.85) \\
Employing two types of feedback and two styles of learning in a ubiquitous learning environment for developing technological skills and technology acceptance among educational technology students & 2024 & OA & supports the claim (0.65) \\
Enhancing Digital Reading Comprehension in Secondary Education: Effects of Question Placement and Feedback Type & 2026 & OA & supports the claim (0.60) \\
Enhancing digital reading comprehension through feedback messages: A large and long-term dynamic approach with secondary school students & 2025 & OA & supports the claim (0.80) \\
Entwicklung und Evaluation von fehlerspezifischem informativem tutoriellem Feedback (ITF) für die schriftliche Subtraktion & 2004 & OA & supports the claim (0.70) \\
Evaluating Feedback Strategies for Virtual Human Trainers & 2020 & OA & supports the claim (0.50) \\
Evaluating the effects of feedback type in a computer-managed learning program & 2021 & OA & supports the claim (0.65) \\
Evaluation of Adaptive Feedback in a Smartphone-Based Game on Health Care Providers' Learning Gain: Randomized Controlled Trial & 2020 & OA & supports the claim (0.70) \\
Everything right or nothing wrong? Regulatory fit effects in an e-learning context & 2022 & OA & supports the claim (0.60) \\
Evidence to support integrating feedback best practice for computerbased assessment in pharmacology education & 2025 & OA & supports the claim (0.90) \\
Explanations and Interactives Improve Subjective Experiences in Online Courseware & 2017 & OA & supports the claim (0.60) \\
EXPLICIT FEEDBACK WITHIN GAME-BASED TRAINING: EXAMINING THE INFLUENCE OF SOURCE MODALITY EFFECTS ON INTERACTION & 2013 & OA & supports the claim (0.35) \\
Exploring how children with reading difficulties respond to instructional supports in literacy games and the role of prior knowledge & 2023 & OA & supports the claim (0.50) \\
Exploring the effects of elaborated and motivational feedback on learning engagement in online scripted role discussion & 2025 & OA & supports the claim (0.50) \\
Factors affecting learning of vector math from computer-based practice: Feedback complexity and prior knowledge & 2016 & OA & supports the claim (0.90) \\
Feedback complexities interact with prior knowledge & 2026 & OA & supports the claim (0.85) \\
Feedback Effects in the Acquisition of a Hierarchical Skill & 2019 & OA & supports the claim (0.60) \\
Feedback for Metacognitive Support in Learning by Teaching Environments & 2006 & OA & supports the claim (0.55) \\
Feedback in Immersive Virtual Reality Learning Environments -- A Scoping Review & 2025 & OA & supports the claim (0.45) \\
Feedback in mathematics education research: a systematic literature review & 2024 & OA & supports the claim (0.55) \\
Feedback That Leads to Improvement in Student Essays: Testing the Hypothesis that ``Where to Next'' Feedback is Most Powerful & 2021 & OA & supports the claim (0.60) \\
Feedback type matters in children's word learning & 2025 & OA & supports the claim (0.60) \\
Feedback, Metacognitive Judgement and Task Difficulty: Determinants of Cognitive Task Performance & 2026 & OA & supports the claim (0.60) \\
Feedback-induced Improvement in the Learning of Anatomy and Physiology-specific Terminology among Tenth Grade High School Students & 2023 & OA & supports the claim (0.65) \\
Feedback-type interactive features in electronic storybooks enhance learning regardless of cognitive differences & 2026 & OA & supports the claim (0.50) \\
Fine-Grained Assessment of Children's Text Comprehension Skills & 2019 & OA & supports the claim (0.45) \\
Formative assessment and feedback for learning in higher education: A systematic review & 2021 & OA & supports the claim (0.70) \\
Formative feedback in online quizzes and its impact on nursing students' self-regulated learning & 2024 & OA & supports the claim (0.50) \\
Fostering achievement and motivation with bug-related tutoring feedback in a computer-based training for written subtraction & 2006 & OA & supports the claim (0.85) \\
Förderung des computerbasierten Wissenserwerbs im Bereich empirischer Forschungsmethoden durch kooperatives Lernen und eine Feedbackmaßnahme & 2003 & OA & supports the claim (0.80) \\
How automated feedback from a digital mathematics textbook affects primary students' conceptual development: two case studies & 2021 & OA & supports the claim (0.50) \\
How do older adults correct memory errors? The effects of practice and metacognitive strategies & 2025 & OA & supports the claim (0.55) \\
How Instructional Feedback Has Been Employed in Instructional Units for Teaching Software Project Management Tools: A Systematic Literature Review & 2016 & OA & supports the claim (0.40) \\
How to Make Computer-Based Feedback More Productive: The Power of Erroneous Solutions & 2024 & OA & supports the claim (0.85) \\
If learners have a choice! A systematic literature review on feedback choice in computer-based learning & 2026 & OA & supports the claim (0.60) \\
Immediate feedback in gamified online Chinese vocabulary quizzing: Effects on perceived cognitive load and strategic vocabulary learning & 2026 & OA & supports the claim (0.40) \\
Impact of Hint Content on Performance and Learning: A Study with Primary School Children in a Scratch Course & 2023 & OA & supports the claim (0.60) \\
Improving Undergraduates' and Postgraduates' Academic Writing Skills with Strategy Training and Feedback & 2017 & OA & supports the claim (0.65) \\
Increased Retest Scores on Cognitive Tests: Learning or Memory Effects? & 2015 & OA & supports the claim (0.45) \\
Increasing Online Information Retention: Analyzing the Effects & 2018 & OA & supports the claim (0.70) \\
Informative tutoring feedback strategies in university calculus teaching: effects on students' performance, self-efficacy and calibration & 2026 & OA & supports the claim (0.60) \\
Inquiry in oral communication: adapting oral examinations for teaching introductory physiology students to evaluate scientific research studies & 2020 & OA & supports the claim (0.60) \\
Intelligent Feedback on Hypothesis Testing & 2020 & OA & supports the claim (0.85) \\
INTERACTIVE EFFECTS OF FEEDBACK TYPE AND FEEDBACK PROPENSITIES ON TASK PERFORMANCE & 2006 & OA & supports the claim (0.60) \\
Interactive Feedback for Learning Mathematics in a Digital Learning Environment & 2021 & OA & supports the claim (0.70) \\
Investigation of feedback on student performance & 2005 & OA & supports the claim (0.80) \\
La importancia del feedback en el aula invertida: motivación y rendimiento académico en universitarios & 2025 & OA & supports the claim (0.55) \\
Learners' use of domain-specific computer-based feedback to overcome logical circularity in deductive proving in geometry & 2018 & OA & supports the claim (0.50) \\
Learning Cognitive Feedback Specificity during Training and the Effect on Learning for Cognitive Tasks & 2009 & OA & supports the claim (0.45) \\
Learning programming online: Influences of various types of feedback on programming performances & 2020 & OA & supports the claim (0.60) \\
Learning to read in first grade through a reading tutor that employs automatic speech recognition: The effects of implicit and explicit feedback & 2025 & OA & supports the claim (0.55) \\
Learning UI Functional Design Principles Through Simulation With Feedback & 2020 & OA & supports the claim (0.45) \\
Levels of exit slip feedback in secondary mathematics & 2017 & OA & supports the claim (0.60) \\
Modeling Alphabet Skills as Instructive Feedback within a Phonological Awareness Curriculum & 2015 & OA & supports the claim (0.40) \\
More Bang for your Buck: Bolstering Learning Via Refutation Text with Refutation-based Elaborated Feedback & 2017 & OA & supports the claim (0.75) \\
O Provimento de Feedback Explanatório em um Objeto de Aprendizagem Multimodal & 2013 & OA & supports the claim (0.50) \\
Online Quizzes with Closed Questions in Formal Assessment: How Elaborate Feedback can Promote Learning & 2020 & OA & supports the claim (0.95) \\
Online self-assessments and student engagement: a causal analysis in university mathematics & 2025 & OA & supports the claim (0.55) \\
Optimizing GAI-assisted formative feedback: an experimental study on its effects on engagement, shared metacognition, and learning performance in online collaborative learning & 2025 & OA & supports the claim (0.70) \\
Personalized feedback in digital learning environments: Classification framework and literature review & 2022 & OA & supports the claim (0.60) \\
Personalized Multimodal Feedback Using Multiple External Representations: Strategy Profiles and Learning in High School Physics & 2026 & OA & supports the claim (0.80) \\
Pilot Trainees Benefit from Modelling and Adaptive Feedback & 2025 & OA & supports the claim (0.55) \\
Predicting feedback effects from response-certitude estimates & 1992 & OA & supports the claim (0.60) \\
Prior knowledge interacts with the effects of prequestions and feedback types on learning from videos: Eyetracking and cognitive load evidence & 2026 & OA & supports the claim (0.70) \\
Program in developing jurisprudential concepts among secondary school students & 2022 & OA & supports the claim (0.60) \\
Providing feedback on computer-based algebra homework in middle-school classrooms & 2016 & OA & supports the claim (0.55) \\
Question-answering skills: The role of feedback in digital environments & 2020 & WoS, OA & supports the claim (0.70) \\
Radiographical diagnostic competences of dental students using various feedback methods and integrating an artificial intelligence application---A randomized clinical trial & 2024 & OA & supports the claim (0.60) \\
Real-time cognitive-affective dynamics of failure feedback in a technology-based learning task & 2026 & OA & supports the claim (0.55) \\
Refining the Retrieval Effort Theory of the Testing Effect as it Relates to Feedback & 2014 & OA & supports the claim (0.50) \\
Repairing Errors With Elaborative Feedback in Computerised Learning Environments & 2026 & OA & supports the claim (0.65) \\
Right or wrong? How feedback content and source influence children's mathematics performance and persistence & 2024 & OA & supports the claim (0.60) \\
Role Of Feedbacks In Simulation-Based Learning & 2016 & OA & supports the claim (0.60) \\
Role of Working Memory and Strategy-Use in Feedback Effects on children's Progression in Analogy Solving:an Explanatory Item Response Theory Account & 2016 & OA & supports the claim (0.85) \\
Rückmeldesysteme in interaktiven Lernvideos für angehende Biologielehrkräfte. Einfluss auf die Forschungskompetenz, das forschungsmethodische Wissen und die intrinsische Motivation & 2024 & OA & supports the claim (0.70) \\
Selective use of animation and feedback in computer-based instruction & 1992 & OA & supports the claim (0.45) \\
Steigerung selbstregulierten Lernens durch computerbasiertes Feedback beim Erwerb von Experimentierkompetenz im Fach Biologie & 2013 & OA & supports the claim (0.55) \\
Students' mathematical feedback processes using a digital curriculum resource: a qualitative perspective & 2026 & OA & supports the claim (0.55) \\
Students' perceptions and outcome of teacher feedback: a systematic review & 2025 & OA & supports the claim (0.50) \\
Supporting learning from text: A meta-analysis on the timing and content of effective feedback & 2019 & OA & supports the claim (0.90) \\
Tailoring Feedback in Online Assessment: Influence of Learning Styles on the Feedback Preferences and Elaborated Feedback Effectiveness & 2008 & IEEE & supports the claim (0.70) \\
Teachers' Feedback on Oral Reading: A Critical Review of its Effects and the use of Theory in Research & 2024 & OA & supports the claim (0.50) \\
Testing pays off twice: Potentials of practice tests and feedback regarding exam performance and judgment accuracy & 2022 & OA & supports the claim (0.50) \\
The benefits of computer-generated feedback for mathematics problem solving & 2016 & OA & supports the claim (0.75) \\
The Development and Evaluation of an Online Formative Assessment upon Single-player Game in E-learning Environment & 2013 & OA & supports the claim (0.80) \\
The Differential Effects of Elaborated Task and Process Feedback on Multi-Digit Multiplication & 2020 & OA & supports the claim (0.50) \\
The effect of feedback on academic achievement in technology-rich learning environments (TREs): A meta-analytic review & 2023 & WoS & supports the claim (0.90) \\
The Effect of Feedback Types on Academic Achievement That Used in Computer Based Training Software & 2011 & OA & supports the claim (0.90) \\
The Effect of Formative Assessments on Teaching and Learning & 2010 & OA & supports the claim (0.60) \\
The effect of immediate feedback and after-action reviews (AARs) on learning, retention and transfer & 2005 & OA & supports the claim (0.45) \\
The effect of negative feedback on motivation : a meta-analytic investigation & 2014 & OA & supports the claim (0.60) \\
The effect of using feedback on developing spelling skills among female students enrolled the Administrative Sciences Diploma program & 2024 & OA & supports the claim (0.50) \\
The Effectiveness of Emotional Motivational Feedback Messages & 2014 & OA & supports the claim (0.45) \\
The effectiveness of explanatory adaptive feedback within a digital educational game to enhance fraction understanding & 2024 & OA & supports the claim (0.75) \\
The effectiveness of methods for providing written feedback through a computer-based assessment for learning: a systematic review & 2011 & WoS & supports the claim (0.80) \\
The effectiveness of methods for providing written feedback through a computer-based assessment for learning: A systematic review & 2011 & Scopus, OA & supports the claim (0.70) \\
The effects of a digital formative assessment tool on mathematics achievement and student motivation: Results of a randomized experiment & 2016 & OA & supports the claim (0.45) \\
The effects of a digital formative assessment tool on spelling achievement: Results of a randomized experiment & 2018 & OA & supports the claim (0.45) \\
The effects of computer-assisted adaptive instruction and elaborated feedback on learning outcomes. A randomized control trial & 2020 & OA & supports the claim (0.70) \\
The Effects of Feedback During Exploration Depend on Prior Knowledge. & 2012 & OA & supports the claim (0.80) \\
The effects of formative feedback on student learning in science education & 2019 & OA & supports the claim (0.60) \\
The effects of instant lab feedback through the use of excel spreadsheet on student learning in high school physics & 2016 & OA & supports the claim (0.50) \\
The effects of process-oriented feedback and revisions on middle schools' fraction arithmetic performance & 2025 & OA & supports the claim (0.80) \\
The electronic formative assessment based on the two types of feedback (interpretive and corrective) and its effectiveness on academic achievement and test anxiety among high school female students in mathematics & 2025 & OA & supports the claim (0.80) \\
The Feedback Principle in Multimedia Learning & 2014 & OA & supports the claim (0.40) \\
The Impact of Computer-Based Feedback Strategies on Students' Learning: A Network Meta-Analysis & 2024 & OA & supports the claim (0.90) \\
The Impact of Emotional Feedback and Elaborated Feedback of a Pedagogical Agent on Multimedia Learning & 2022 & OA & supports the claim (0.90) \\
The Impact of Post-Retrieval Practice and Feedback Types on the Accuracy of Metacognitive Judgments and Academic Achievement in Online Learning & 2024 & OA & supports the claim (0.80) \\
The Impact of Prompts and Feedback on the Performance during Multi-Session Self-Regulated Learning in the Hypermedia Environment & 2023 & OA & supports the claim (0.60) \\
The Influence of Feedback Content and Feedback Time on Multimedia Learning Achievement of College Students and Its Mechanism & 2021 & OA & supports the claim (0.80) \\
The interaction effect of pedagogical agent and emotional feedback on effective learning: a 2 × 2 factorial experiment in online formative assessment & 2025 & OA & supports the claim (0.50) \\
The more, the better? Learning with feedback and instruction & 2023 & OA & supports the claim (0.50) \\
The Next Black Box of Formative Assessment: A Model of the Internal Mechanisms of Feedback Processing & 2022 & OA & supports the claim (0.40) \\
The role of feedback type and peer interaction on knowledge acquisition in a flipped classroom on social science research methods & 2024 & OA & supports the claim (0.85) \\
The Type and Timing of Feedback within an Intelligent Console-Operations Tutor & 1994 & OA & supports the claim (0.80) \\
The use of feedback in web-based instruction : achievement, feedback study time, and efficiency & 2009 & OA & supports the claim (0.55) \\
Time is of the Essence: Log-Based Investigation of the Interplay between Feedback Elaboration, Time to Process Feedback, and Task Success in Online Mathematics & 2023 & OA & supports the claim (0.75) \\
Toward self-regulated learning: effects of different types of data-driven feedback on pupils' mathematics word problem-solving performance & 2024 & OA & supports the claim (0.60) \\
Towards reducing anxiety and increasing performance in physics education: Evidence from a randomized experiment. & 2019 & OA & supports the claim (0.50) \\
Type and timing of feedback within an intelligent console-operations tutor & 1994 & Scopus & supports the claim (0.40) \\
Understanding individual differences in students' responses to technology-based feedback on a writing task: the role of achievement motives and initial task performance & 2025 & OA & supports the claim (0.60) \\
Understanding the Hypercorrection Effect: Why High-Confidence Errors are More Likely to be Corrected & 2010 & OA & supports the claim (0.55) \\
Understanding variation in children's reading comprehension: A dynamic approach & 2023 & OA & supports the claim (0.50) \\
Unfolding Learners' Response to Different Versions of Automated Feedback in a MOOC for Programming -- A Sequence Analysis Approach & 2023 & OA & supports the claim (0.60) \\
Use of practice tests with immediate feedback in an undergraduate molecular biology course & 2022 & OA & supports the claim (0.55) \\
Using digital formative assessments to improve learning in physics education & 2021 & OA & supports the claim (0.55) \\
Utility of Feedback Has a Greater Impact on Learning than Ease of Decoding & 2019 & OA & supports the claim (0.50) \\
Uyarlanabilir Dönüt Sistemi Tasarımı İçin Kullanıcı Profillerinin Belirlenmesi & 2019 & OA & supports the claim (0.50) \\
What makes tabletbased learning effective? A study of the role of realtime adaptive feedback & 2024 & OA & supports the claim (0.70) \\
When to Quiz in Instructional Videos? Effects on Learning and SelfRegulated Learning in Secondary Education & 2026 & OA & supports the claim (0.75) \\
(De)motivating ZeroPerforming Students With Negative Feedback: Does the Salience of Performance Information Matter? & 2025 & OA & qualifies or contradicts the claim (0.55) \\
``I just looked at the answers'': How the voluntary use of online quizzes impacts performance in introductory accounting. & 2021 & OA & qualifies or contradicts the claim (0.65) \\
Accuracy Matters for the Benefits of Sleep After Retrieval Practice & 2019 & OA & qualifies or contradicts the claim (0.55) \\
Adaptive training instructional interventions: A meta-analysis & 2024 & OA & qualifies or contradicts the claim (0.50) \\
Am I Wrong, or Is the Autograder Wrong? Effects of AI Grading Mistakes on Learning & 2023 & OA & qualifies or contradicts the claim (0.70) \\
Beyond access: Engagement with conversational AI-assisted feedback shapes cognitive load in digital game-based learning & 2026 & OA & qualifies or contradicts the claim (0.50) \\
Case study: a comparison of immediate and delayed feedback, in the context of online testing with fourth class students & 2013 & OA & qualifies or contradicts the claim (0.70) \\
Channelling feedback through audiovisual presentations: Do higher education students perceive, use and benefit from video feedback compared to written feedback? & 2024 & OA & qualifies or contradicts the claim (0.75) \\
Combined inner and outer loop feedback in an intelligent tutoring system for statistics in higher education & 2020 & OA & qualifies or contradicts the claim (0.85) \\
Computer-Based Formative Assessment : Variables Influencing Feedback Behaviour & 2013 & OA & qualifies or contradicts the claim (0.70) \\
Content and timing of feedback in a web-based learning environment: effects on learning as a function of prior knowledge & 2008 & OA & qualifies or contradicts the claim (0.80) \\
Delayed Versus Immediate Feedback in an Independent Study High School Setting & 2005 & OA & qualifies or contradicts the claim (0.55) \\
Delaying elaborated feedback within computerbased learning environments: The role of summative and questionbased feedback & 2021 & OA & qualifies or contradicts the claim (0.90) \\
Difficulties of adults with low educational level in using computerized explanatory feedback & 1997 & Scopus & qualifies or contradicts the claim (0.50) \\
Do differences in topic knowledge matter? An experimental investigation into topic knowledge as a possible moderator of the testing effect & 2025 & OA & qualifies or contradicts the claim (0.50) \\
Does providing the correct diagnosis as feedback after self-explanation improve medical students diagnostic performance? & 2019 & OA & qualifies or contradicts the claim (0.70) \\
Effects of AI generated adaptive feedback on statistical skills and interest in statistics: A field experiment in higher education & 2025 & OA & qualifies or contradicts the claim (0.85) \\
Effects of elaborate feedback during practice tests: Costs and benefits of retrieval prompts. & 2019 & OA & qualifies or contradicts the claim (0.90) \\
Effects of Elaborative Feedback and Correct Answer Feedback on Reducing Confirmation Bias: A Preregistered Study & 2019 & OA & qualifies or contradicts the claim (0.85) \\
EFL Learners Interaction with Feedback Presented through a Computer-Assisted Reading Program & 2022 & OA & qualifies or contradicts the claim (0.80) \\
Elaborative Feedback to Enhance Online Second Language Reading Comprehension & 2017 & OA & qualifies or contradicts the claim (0.90) \\
Examining the Role of Feedback on Self-efficacy and Performance in Web-based Environment & 2007 & IEEE & qualifies or contradicts the claim (0.70) \\
Feedback Timing: The Impact on Undergraduates' Statistics Anxiety, Self-Efficacy and Perceptions & 2026 & OA & qualifies or contradicts the claim (0.85) \\
Formative Assessment in Mathematics Education: A Systematic Review & 2025 & OA & qualifies or contradicts the claim (0.50) \\
Formative Assessment: Comparing Immediate and Delayed Feedback & 2012 & OA & qualifies or contradicts the claim (0.70) \\
Fostering pre-service teachers' classroom management knowledge, self-efficacy, and professional vision: The effect of different expert feedback and pre-service teachers' feedback perceptions during online video analysis & 2025 & OA & qualifies or contradicts the claim (0.70) \\
From belief to evidence: simpler immediate feedback improves language learning and confidence in semi-open-ended questions & 2025 & OA & qualifies or contradicts the claim (0.90) \\
How do students process complex formative feedback in question-answering tasks? A think-aloud study & 2019 & OA & qualifies or contradicts the claim (0.75) \\
How to Make Computer-Based Feedback More Effective: The Benefits of Allowing Learners to Choose Feedback Types & 2026 & WoS & qualifies or contradicts the claim (0.60) \\
Impact of different practice testing methods on learning outcomes & 2024 & OA & qualifies or contradicts the claim (0.60) \\
Impact of feedback type on pre-service teachers' professional vision, self-efficacy, and knowledge & 2026 & OA & qualifies or contradicts the claim (0.60) \\
Instructional advice, time advice and learning questions in computer simulations & 2010 & OA & qualifies or contradicts the claim (0.40) \\
Investigating a blended learning context that incorporates two-stage quizzes and peer formative feedback in STEM education & 2022 & OA & qualifies or contradicts the claim (0.45) \\
Is technology always helpful?: A critical review of the impact on learning outcomes of education technology in supporting formative assessment in schools & 2021 & OA & qualifies or contradicts the claim (0.60) \\
LEARNING ANATOMY WITH AUGMENTED REALITY & 2018 & OA & qualifies or contradicts the claim (0.40) \\
Learning and Motivation When Using Multiple-Try in a Digital Game for Primary Students in Chile & 2023 & OA & qualifies or contradicts the claim (0.60) \\
Learning to read and understand philosophy texts through feedback: The effectiveness of different feedback variants in the ill-structured learning domain of the philosophy of science & 2026 & OA & qualifies or contradicts the claim (0.70) \\
Lernaufgaben mit automatisiertem Feedback in einem digitalen Hochschultutorium zur Vorlesung der Allgemeinen Chemie als Angebot zur Binnendifferenzierung & 2023 & OA & qualifies or contradicts the claim (0.80) \\
LLM-Generated Feedback Supports Learning if Learners Choose to Use it & 2025 & OA & qualifies or contradicts the claim (0.70) \\
Maximizing feedback for language learning: English language learners' attention, affect, cognition and usage of computer-delivered feedback from an English language reading proficiency assessment & 2017 & OA & qualifies or contradicts the claim (0.55) \\
Older Adults With Cognitive Impairment Can Correct COVID19 HealthRelated Misconceptions & 2025 & OA & qualifies or contradicts the claim (0.40) \\
Reading comprehension exercises online: The effects of feedback, proficiency and interaction & 2007 & WoS, Scopus, OA & qualifies or contradicts the claim (0.50) \\
Retrieval Practice in Classroom Settings: A Review of Applied Research & 2019 & OA & qualifies or contradicts the claim (0.55) \\
Rubrics and in-class feedback on a higher education course with statistical content & 2019 & OA & qualifies or contradicts the claim (0.55) \\
Seeking out fun failure: how positive failure feedback could enhance the instructional effectiveness of CALL mini-games & 2013 & OA & qualifies or contradicts the claim (0.50) \\
Self-Regulation and Feedback in an Educational Statistics Course & 2018 & OA & qualifies or contradicts the claim (0.80) \\
Student Preferences for Multiple Attempts and Feedback on Online Quantitative Assessments & 2023 & OA & qualifies or contradicts the claim (0.50) \\
Students' Verification and Elaboration in Outdoor Mathematics: The Role of Digital Feedback in MathCityMap & 2024 & OA & qualifies or contradicts the claim (0.70) \\
Supporting comprehension: The advantages of multiple--choice over true---false practice tests & 2025 & OA & qualifies or contradicts the claim (0.70) \\
Testing the boundaries of the hypercorrection effect: when metacognitive accuracy facilitates and undermines learning in adults & 2026 & OA & qualifies or contradicts the claim (0.60) \\
The effect of feedback complexity on university students' achievement and intrinsic motivation & 2020 & OA & qualifies or contradicts the claim (0.85) \\
The effect of the frequency of use of an intelligent tutoring system on learning gains in mathematics secondary education & 2026 & OA & qualifies or contradicts the claim (0.50) \\
The effect of the nonhuman external regulator's answer-until-correct (AUC) versus knowledge-of-result (KR) task feedback on children's behavioral regulation during learning tasks & 2011 & OA & qualifies or contradicts the claim (0.40) \\
The effect of tutorial feedback type on the choice of feedback type in pre-service teachers' development of learning objects & 2017 & OA & qualifies or contradicts the claim (0.35) \\
The Effect of Whole and Part-Task Training and Feedback during Simulated Instrument Flight Training & 2009 & OA & qualifies or contradicts the claim (0.45) \\
The Effects of Different Levels of Performance Feedback on TOEFL iBT ® Reading Practice Test Performance & 2017 & OA & qualifies or contradicts the claim (0.92) \\
THE EFFECTS OF IMMEDIATE AND DELAYED FEEDBACK ON SECONDARY DISTANCE LEARNERS & 2007 & OA & qualifies or contradicts the claim (0.55) \\
The Effects of Item-by-Item Feedback during a Computer Adaptive Test. & 2014 & OA & qualifies or contradicts the claim (0.50) \\
The effects of various conditions of adjunctive interactive computer-assisted testing on final examination performance / & 1978 & OA & qualifies or contradicts the claim (0.50) \\
The Impact of Volitional Feedback on Learners' Self-Efficacy and Course Satisfaction in a College Assignment System & 2011 & OA & qualifies or contradicts the claim (0.50) \\
The influence of evaluative right/wrong feedback on phonological and semantic processes in word learning & 2018 & OA & qualifies or contradicts the claim (0.75) \\
The interplay of task difficulty, precision and feedback in students' help-seeking behaviour in online maths environments & 2026 & OA & qualifies or contradicts the claim (0.55) \\
The Preparation and Use of Self-Access English CALL Courses for Junior High School Students & 2006 & OA & qualifies or contradicts the claim (0.40) \\
The relationship between formative assessment and reading achievement: A multilevel analysis of students in 19 countries/regions & 2022 & OA & qualifies or contradicts the claim (0.45) \\
The role of teacher feedback on students' motivation: a situated expectancy-value perspective & 2026 & OA & qualifies or contradicts the claim (0.55) \\
Three forms of feedback in digital learning material in mathematics & 2025 & OA & qualifies or contradicts the claim (0.90) \\
Time advice and learning questions in computer simulations & 2011 & OA & qualifies or contradicts the claim (0.70) \\
Timing is everything: investigating users' feedback timing and content preferences to enhance computer-based training & 2026 & OA & qualifies or contradicts the claim (0.70) \\
Tools that detectives use: in search of learner-related determinants for usage of optional feedback in a written murder mystery & 2013 & OA & qualifies or contradicts the claim (0.55) \\
Understanding student behaviors using immediate feedback features in a blended learning environment & 2015 & OA & qualifies or contradicts the claim (0.55) \\
University Students' Engagement with Feedback in a Quiz Platform: A Case Study of Quizizz & 2024 & OA & qualifies or contradicts the claim (0.55) \\
Unlocking math potential in students from lower SES backgrounds -- using instructional scaffolds to improve performance & 2025 & OA & qualifies or contradicts the claim (0.60) \\
Using eye-tracking to measure cognitive engagement with feedback in a digital literacy game & 2023 & OA & qualifies or contradicts the claim (0.60) \\
Using Sequential Pattern Mining to Understand How Students Use Guidance While Doing Scientific Calculations & 2023 & OA & qualifies or contradicts the claim (0.55) \\
Weak memory for future-oriented feedback: investigating the roles of attention and improvement focus & 2019 & OA & qualifies or contradicts the claim (0.50) \\
What's Wrong With Giving Students Feedback? & 2020 & OA & qualifies or contradicts the claim (0.80) \\
Where less is more: Limited feedback in formative online multiplechoice tests improves student selfregulation & 2023 & OA & qualifies or contradicts the claim (0.85) \\
\end{longtable}
\endgroup

